% Options for packages loaded elsewhere
\PassOptionsToPackage{unicode}{hyperref}
\PassOptionsToPackage{hyphens}{url}
\documentclass[
  sigconf,
  anonymous]{acmart}
\usepackage{xcolor}
\usepackage{amsmath}
\setcounter{secnumdepth}{5}
\usepackage{iftex}
\ifPDFTeX
  \usepackage[T1]{fontenc}
  \usepackage[utf8]{inputenc}
  \usepackage{textcomp} % provide euro and other symbols
\else % if luatex or xetex
  \usepackage{unicode-math} % this also loads fontspec
  \defaultfontfeatures{Scale=MatchLowercase}
  \defaultfontfeatures[\rmfamily]{Ligatures=TeX,Scale=1}
\fi
\ifPDFTeX\else
  % xetex/luatex font selection
\fi
% Use upquote if available, for straight quotes in verbatim environments
\IfFileExists{upquote.sty}{\usepackage{upquote}}{}
\IfFileExists{microtype.sty}{% use microtype if available
  \usepackage[]{microtype}
  \UseMicrotypeSet[protrusion]{basicmath} % disable protrusion for tt fonts
}{}
\makeatletter
\@ifundefined{KOMAClassName}{% if non-KOMA class
  \IfFileExists{parskip.sty}{%
    \usepackage{parskip}
  }{% else
    \setlength{\parindent}{0pt}
    \setlength{\parskip}{6pt plus 2pt minus 1pt}}
}{% if KOMA class
  \KOMAoptions{parskip=half}}
\makeatother
\usepackage{color}
\usepackage{fancyvrb}
\newcommand{\VerbBar}{|}
\newcommand{\VERB}{\Verb[commandchars=\\\{\}]}
\DefineVerbatimEnvironment{Highlighting}{Verbatim}{commandchars=\\\{\}}
% Add ',fontsize=\small' for more characters per line
\newenvironment{Shaded}{}{}
\newcommand{\AlertTok}[1]{\textcolor[rgb]{1.00,0.00,0.00}{\textbf{#1}}}
\newcommand{\AnnotationTok}[1]{\textcolor[rgb]{0.38,0.63,0.69}{\textbf{\textit{#1}}}}
\newcommand{\AttributeTok}[1]{\textcolor[rgb]{0.49,0.56,0.16}{#1}}
\newcommand{\BaseNTok}[1]{\textcolor[rgb]{0.25,0.63,0.44}{#1}}
\newcommand{\BuiltInTok}[1]{\textcolor[rgb]{0.00,0.50,0.00}{#1}}
\newcommand{\CharTok}[1]{\textcolor[rgb]{0.25,0.44,0.63}{#1}}
\newcommand{\CommentTok}[1]{\textcolor[rgb]{0.38,0.63,0.69}{\textit{#1}}}
\newcommand{\CommentVarTok}[1]{\textcolor[rgb]{0.38,0.63,0.69}{\textbf{\textit{#1}}}}
\newcommand{\ConstantTok}[1]{\textcolor[rgb]{0.53,0.00,0.00}{#1}}
\newcommand{\ControlFlowTok}[1]{\textcolor[rgb]{0.00,0.44,0.13}{\textbf{#1}}}
\newcommand{\DataTypeTok}[1]{\textcolor[rgb]{0.56,0.13,0.00}{#1}}
\newcommand{\DecValTok}[1]{\textcolor[rgb]{0.25,0.63,0.44}{#1}}
\newcommand{\DocumentationTok}[1]{\textcolor[rgb]{0.73,0.13,0.13}{\textit{#1}}}
\newcommand{\ErrorTok}[1]{\textcolor[rgb]{1.00,0.00,0.00}{\textbf{#1}}}
\newcommand{\ExtensionTok}[1]{#1}
\newcommand{\FloatTok}[1]{\textcolor[rgb]{0.25,0.63,0.44}{#1}}
\newcommand{\FunctionTok}[1]{\textcolor[rgb]{0.02,0.16,0.49}{#1}}
\newcommand{\ImportTok}[1]{\textcolor[rgb]{0.00,0.50,0.00}{\textbf{#1}}}
\newcommand{\InformationTok}[1]{\textcolor[rgb]{0.38,0.63,0.69}{\textbf{\textit{#1}}}}
\newcommand{\KeywordTok}[1]{\textcolor[rgb]{0.00,0.44,0.13}{\textbf{#1}}}
\newcommand{\NormalTok}[1]{#1}
\newcommand{\OperatorTok}[1]{\textcolor[rgb]{0.40,0.40,0.40}{#1}}
\newcommand{\OtherTok}[1]{\textcolor[rgb]{0.00,0.44,0.13}{#1}}
\newcommand{\PreprocessorTok}[1]{\textcolor[rgb]{0.74,0.48,0.00}{#1}}
\newcommand{\RegionMarkerTok}[1]{#1}
\newcommand{\SpecialCharTok}[1]{\textcolor[rgb]{0.25,0.44,0.63}{#1}}
\newcommand{\SpecialStringTok}[1]{\textcolor[rgb]{0.73,0.40,0.53}{#1}}
\newcommand{\StringTok}[1]{\textcolor[rgb]{0.25,0.44,0.63}{#1}}
\newcommand{\VariableTok}[1]{\textcolor[rgb]{0.10,0.09,0.49}{#1}}
\newcommand{\VerbatimStringTok}[1]{\textcolor[rgb]{0.25,0.44,0.63}{#1}}
\newcommand{\WarningTok}[1]{\textcolor[rgb]{0.38,0.63,0.69}{\textbf{\textit{#1}}}}
\usepackage{longtable,booktabs,array}
\newcounter{none} % for unnumbered tables
\usepackage{calc} % for calculating minipage widths
% Correct order of tables after \paragraph or \subparagraph
\usepackage{etoolbox}
\makeatletter
\patchcmd\longtable{\par}{\if@noskipsec\mbox{}\fi\par}{}{}
\makeatother
% Allow footnotes in longtable head/foot
\IfFileExists{footnotehyper.sty}{\usepackage{footnotehyper}}{\usepackage{footnote}}
\makesavenoteenv{longtable}
\setlength{\emergencystretch}{3em} % prevent overfull lines
\providecommand{\tightlist}{%
  \setlength{\itemsep}{0pt}\setlength{\parskip}{0pt}}
\setmonofont{DejaVuSansMono.ttf}[Scale=0.78]
\usepackage[htt]{hyphenat}
\usepackage{fvextra}
\RecustomVerbatimEnvironment{verbatim}{Verbatim}{fontsize=\scriptsize,breaklines,breakanywhere}
\emergencystretch=3em
% Math macros used by the paper that standard LaTeX lacks:
\DeclareMathOperator*{\argmin}{arg\,min}
\newcommand{\llbracket}{[\![}
\newcommand{\rrbracket}{]\!]}
\usepackage{bookmark}
\IfFileExists{xurl.sty}{\usepackage{xurl}}{} % add URL line breaks if available
\urlstyle{same}
\hypersetup{
  pdftitle={AgentISA: Formal Instruction Set Design for LLM Agents},
  pdfauthor={Anonymous Author(s)},
  hidelinks,
  pdfcreator={LaTeX via pandoc}}

\title{AgentISA: Formal Instruction Set Design for LLM Agents}
\author{Anonymous Author(s)}
\date{}

\begin{document}
\begin{abstract}
Large Language Models run as autonomous agents, yet ad-hoc prompt
engineering and per-provider function calling leave their behavior
difficult to specify, analyze, and optimize systematically. We present
AgentISA, an instruction set architecture making agent behavior a
first-class compilable artifact: (1) a formal instruction
set---compositional semantics, a type system,
pre/post-conditions---checked offline by a shipped bounded verifier for
pure instructions, bracketed by runtime pre-dispatch validation; (2)
Agentic Calling, lifting provider-specific function calling into
reusable, cacheable instructions; (3) instruction learning that mines
recurring trace patterns, promoting them as composite-replay
instructions; (4) a compiler and runtime applying classical
optimization---fusion, dead-code elimination, pipeline scheduling,
dynamic hot-path specialization---to agent instruction streams.

A reproducible microbenchmark drives the real implementation against a
mock LLM backend: result caching reduces simulated API calls by
59.5--79.6\% on synthetic Zipf workloads (2,000 tasks, 5 seeds); static
compilation cuts the effective instruction count by 34.2\% (35.4\% at
length 1,000), decomposed by a same-seed ablation into fusion (21.6\%)
and dead-code elimination (24.5\%) with an 11.9-point overlap; dynamic
hot-path optimization shrinks repeatedly executed traces by 65.8\%. Two
real-gateway experiments extend validation: the cache reproduces
computed hit/miss statistics at every skew level with measured
wall-clock and token savings, and the compiled, hot-path-adopted stream
issues zero new gateway calls with byte-identical artifacts
(cache-mediated). A four-task system-level characterization
(12.5--25.0\% per-task compile shrink) calibrates the synthetic
headline; 43 multi-agent tests provide functional validation. We
position these results as controlled-setting evidence, not as production
end-to-end measurements. The shipped prototype implements result
caching, the benchmark-firing static passes, and threshold-triggered
trace re-optimization; tree-pattern instruction selection,
cycle-accurate pipeline scheduling, symbolic-equivalence checking, and
capability metric extraction remain specified or stubbed---no claim
rests on them.
\end{abstract}
\ccsdesc[500]{Software and its engineering~\textgreater{}Compilers}
\ccsdesc[500]{Theory of computation~\textgreater{}Program reasoning}
\ccsdesc[500]{Computing methodologies~\textgreater{}Artificial intelligence}
\ccsdesc[500]{Computer systems organization~\textgreater{}Architectures}
\keywords{Instruction Set Architecture; Large Language Models; Compiler Optimization; Multi-agent Systems; Formal Methods; Result Caching}
\maketitle

\section{Introduction}\label{introduction}

Recent advances in Large Language Models (LLMs) have produced systems
with strong understanding, generation, and---through instruction tuning
and human feedback---instruction-following capabilities {[}1, 2, 3, 4,
5, 6{]}, and a rapidly growing body of work couples them with tools,
memory, and control loops to build autonomous agents {[}9, 10, 11,
12{]}. Yet the engineering substrate beneath these agents remains ad
hoc: behaviors are specified as prompts and provider-specific
function-calling schemas---difficult to specify precisely, analyze
formally, or optimize systematically.

\subsection{Background and Motivation}\label{background-and-motivation}

Function calling---selecting an API from a declared tool set and
emitting structured arguments---has become the dominant interface
between LLMs and external systems {[}13, 14, 15{]}, consolidated by
standardization efforts such as the Model Context Protocol {[}18{]} and
measured by harnesses such as the Berkeley Function Calling Leaderboard
{[}17{]}. However, the unit of interaction remains the individual API
call: each call is dispatched, priced, and forgotten independently, even
when an agent repeatedly executes structurally identical work.

This per-call statelessness contrasts with how computing systems
historically handle repetitive execution: classical compilers and
hardware translate recurring patterns into cheaper forms---instruction
fusion, dead-code elimination, caching, trace-based dynamic optimization
{[}33, 34, 45{]}. Our motivating hypothesis is that these techniques
apply to LLM agent workloads: if agent behavior is expressed as an
explicit instruction stream rather than opaque prompts, recurring
patterns become visible to a compiler and optimizable with
well-understood, semantics-preserving transformations.

\subsection{Challenges in Current LLM-based Agent
Systems}\label{challenges-in-current-llm-based-agent-systems}

Five structural challenges motivate AgentISA:

\begin{enumerate}
\def\labelenumi{\arabic{enumi}.}
\item
  \textbf{Per-call statelessness of tool interfaces.} Mainstream
  function-calling protocols treat each call as an isolated event {[}13,
  14, 18{]}: the provider has no notion of a \emph{program} being
  executed, repeated identical work incurs repeated cost, and caching,
  where present, is a bolt-on afterthought rather than a property of the
  execution model.
\item
  \textbf{Experience that does not accumulate as compiler-visible
  instruction structure.} Paradigms such as Reflexion {[}10{]} and
  Voyager {[}12{]} accumulate experience as verbal reflections or
  inference-time skill libraries---not reified as reusable, statically
  optimizable instructions that a compiler can fuse, schedule, or
  eliminate.
\item
  \textbf{Absence of a shared behavioral contract.} Each multi-agent
  framework defines its own abstractions for roles, messages, and
  workflows {[}19, 20, 21, 22, 23, 24{]}; no instruction-level contract
  spans frameworks, fragmenting tooling and cross-framework analysis.
\item
  \textbf{Cost pressure at the wrong layer.} FrugalGPT {[}28{]} shows
  LLM inference cost is a first-order concern, and serving systems
  attack it below the API boundary with paged attention, continuous
  batching, KV-cache-centric architectures, and attention kernels {[}38,
  39, 40, 41{]}; what is missing is the complementary, above-the-API
  layer that reduces the \emph{number and size} of calls an agent issues
  in the first place.
\item
  \textbf{Integration complexity across providers.} Portability is
  handled at the library level {[}26, 27{]}: logic written against one
  provider's function-calling surface does not transfer cleanly to
  another, and capability differences are hard to reason about
  systematically {[}17{]}.
\end{enumerate}

\subsection{Research Objectives and
Contributions}\label{research-objectives-and-contributions}

This paper introduces AgentISA, an instruction set architecture that
addresses these challenges by making agent behavior a first-class
compilable artifact. Our contributions are:

\begin{enumerate}
\def\labelenumi{\arabic{enumi}.}
\item
  \textbf{A formal instruction set} for LLM agents with compositional
  semantics, a type system, and pre/post-conditions on every instruction
  (Section 3.1). Two deliberately modest sanity results are stated as
  design rationale: Turing completeness (Theorem 3.1)---negative
  certification that the instruction set does not constrain
  expressiveness---and type safety for the pure core (Theorem 3.2),
  which fixes the well-formedness obligation the offline bounded checker
  discharges for pure instructions (Section 3.3.3) and the runtime
  brackets with pre-dispatch context validation.
\item
  \textbf{Agentic Calling}, a mechanism lifting provider-specific
  function calling into reusable, cacheable, semantically specified
  agent instructions, with a model-adapter layer for cross-provider
  portability (Section 3.2).
\item
  \textbf{Trace-based instruction learning}, which mines recurring
  instruction patterns from interaction traces by frequency-based
  sequence mining and deterministic behavior grouping, and promotes
  validated recurring patterns into candidate instructions whose
  \texttt{execute} bodies replay the mined constituent subsequence by
  sequential composition over a shared context (Section 3.3; replay
  reproduces the recorded subsequence's outputs on the same inputs and
  does not generalize beyond them, as the section states).
\item
  \textbf{A compiler and runtime} implementing fusion, dead-code
  elimination, parallel-grouped scheduling, and dynamic hot-path
  specialization for agent instruction streams (Sections 3.4 and 4),
  together with specified tree-pattern instruction selection and
  cycle-accurate pipeline scheduling (Sections 4.1.1 and 4.1.2), with
  semantic-preservation arguments for the implemented rewrites (Theorems
  3.4--3.6, proofs in Appendix C).
\item
  \textbf{An empirical evaluation} of the implemented system, with full
  methodology disclosure. On synthetic Zipf-distributed workloads (5
  random seeds each), caching reduces simulated API calls by
  59.5--79.6\% depending on skew, static compilation reduces the
  effective instruction count by 34.2\% \(\pm\) 1.8\% (sequence length
  200), and dynamic optimization shrinks repeatedly executed traces by
  65.8\% \(\pm\) 4.6\%. Two real-gateway experiments (Sections 5.7 and
  5.7.1) validate the cache and the compiled pipeline against a real LLM
  gateway, and a four-task system-level characterization (Section 5.8)
  composes the mechanisms on realistic (non-synthetic) instruction
  streams.
\item
  \textbf{Multi-agent coordination infrastructure} with capability-based
  agent selection, execution planning, fallback strategies, and a
  security monitor enforcing authentication, authorization,
  resource-access, and data-sensitivity policies (Sections 3.5 and 5.6),
  validated by 43 passing functional tests.
\end{enumerate}

Our empirical claims are scoped accordingly: Sections 5.2--5.5 exercise
the real \texttt{agentflow.core.isa} implementation against a mock LLM
backend on synthetic workloads; Sections 5.7--5.8 extend the validation
at reduced scale on realistic streams; none are production end-to-end
measurements (threats to validity in Section 5.9).

Section 2 reviews related work; Sections 3 and 4 present the design and
its implementation; Section 5 the empirical evaluation; Sections 6 and 7
discussion and conclusions.

\section{Related Work}\label{related-work}

AgentISA draws on eight lines of work, reviewed in turn.

\subsection{Foundation LLMs and Function
Calling}\label{foundation-llms-and-function-calling}

Modern LLMs are autoregressive transformers trained at large scale {[}1,
2, 3, 4{]}, with instruction tuning and reinforcement learning from
human feedback making them reliable instruction followers {[}5, 6{]}.
Their text distribution is conventionally factorized as:

\[
P(y|x) = \prod_{t=1}^{T} P(y_t|x, y_{<t}; \theta)
\]

where \(x\) is the input context, \(y\) the generated output, and
\(\theta\) the model parameters. Function calling constrains this
distribution to a structured output space:

\[
F(x) = \text{argmax}_{f \in \mathcal{F}} P(f|x, \mathcal{F})
\]

where \(\mathcal{F}\) is the set of declared functions and \(F(x)\) the
selected call. GPT-4 {[}3{]}, PaLM {[}2{]}, and open models such as
LLaMA {[}4{]} expose this capability with provider-specific schemas;
AgentISA does not modify this layer but treats any such interface as a
\emph{backend} targeted by a common instruction set, as a compiler
targets multiple processor backends {[}33{]}.

The cost structure of this layer motivates our work. If each call
carries latency \(\tau_i\) plus a context-switching overhead
\(\omega_i\) (prompt reconstruction, serialization, cache misses; the
symbol \(\delta\) is reused in Appendix B for per-instruction transition
functions---the two are unrelated), then \(n\) similar calls cost:

\[
C(n) = \sum_{i=1}^{n} (\tau_i + \omega_i)
\]

\textbf{Definition 2.1} (Function Call Overhead): The overhead of
executing \(n\) function calls is \(C(n)\), the sum of per-call
latencies and per-call context reconstruction costs. When repeated calls
carry identical semantic inputs and satisfy the cache-validity
conditions of Section 3.2, the repeated context-construction cost
\(\omega_i\) is avoidable---which is precisely the cost AgentISA's
caching and fusion are designed to remove. Structural similarity alone
does not make a re-invocation waste: if an input's dependencies have
changed or the call probes external state, re-execution has semantic
meaning, and the cache-validity conditions---not structural
identity---are what license reuse.

\subsection{Reasoning and Agent
Paradigms}\label{reasoning-and-agent-paradigms}

A second line enhances reasoning at prompt time: chain-of-thought
{[}7{]} and zero-shot reasoning {[}8{]}; ReAct's interleaved reasoning
and action {[}9{]}; Reflexion's verbal self-critique {[}10{]}; Tree of
Thoughts' deliberative search {[}11{]}; Voyager's open-ended skill
acquisition {[}12{]}. These paradigms define \emph{what} an agent should
do at the cognitive level; AgentISA defines \emph{how} the resulting
behaviors are represented, compiled, and executed as instructions---a
ReAct loop becomes an instruction sequence whose tool invocations are
fused, cacheable agentic calls.

\subsection{Tool Calling and Interface
Standardization}\label{tool-calling-and-interface-standardization}

Toolformer trains models to interleave tool calls with text {[}13{]};
Gorilla connects models to large API corpora {[}14{]}; ToolLLM scales
tool coverage to 16,000+ real-world APIs {[}15{]}; HuggingGPT
orchestrates model hubs {[}16{]}. On the standardization side, the Model
Context Protocol specifies how tools and context are exposed {[}18{]},
and the Berkeley Function Calling Leaderboard evaluates function-calling
accuracy cross-model {[}17{]}. These efforts standardize the
\emph{interface}; AgentISA standardizes the \emph{execution model above
it}: an AgentISA program compiled for two providers differs only in the
final translation step (Section 3.2.3), and interface-level standards
alone cannot fuse or eliminate calls.

\subsection{Multi-Agent Frameworks}\label{multi-agent-frameworks}

AutoGen frames agent applications as multi-agent conversations {[}19{]};
MetaGPT assigns standardized roles and workflows {[}20{]}; AgentVerse
facilitates multi-agent collaboration and studies emergent behaviors
{[}21{]}; ChatDev organizes agents as a software-development
organization {[}22{]}; CAMEL studies communicative agent societies
{[}23{]}; AutoGPT popularized autonomous goal-driven agents {[}24{]}.
AgentISA contributes a different artifact to this space: not a framework
or orchestration philosophy but an instruction-level
substrate---capability profiles, coordination instructions with barrier
semantics, and a security monitor providing capability-based agent
selection, per-stage fallback planning, and policy enforcement for
authentication, authorization, resource access, and data
sensitivity---that any such framework could target (Section 3.5, Section
5.6).

\subsection{Prompt Compilation and Cost
Optimization}\label{prompt-compilation-and-cost-optimization}

Closest to our compilation perspective is DSPy, which compiles
declarative prompt programs into optimized pipelines via prompt search
and fine-tuning {[}25{]}; LangChain {[}26{]} and LlamaIndex {[}27{]}
provide library-level chaining and indexing abstractions; FrugalGPT
reduces cost through model cascades {[}28{]}. AgentISA differs in the
level at which optimization operates. DSPy optimizes \emph{prompts};
FrugalGPT optimizes \emph{model selection}; AgentISA optimizes
\emph{instruction streams}: fusion, dead-code elimination, parallel
grouping (the shipped form of pipeline scheduling), result caching, and
dynamic hot-path specialization---implemented and measured (Sections
3.4, 4, and 5)---plus tree-pattern instruction selection over pattern
covers and cycle-accurate pipeline scheduling, specified but not yet
implemented (Sections 4.1.1 and 4.1.2). These classical compiler
transformations {[}33{]} compose with---rather than
replace---prompt-level and model-level optimization.

\subsection{Benchmarks for Agent
Systems}\label{benchmarks-for-agent-systems}

AgentBench evaluates LLMs as agents across diverse environments
{[}29{]}; GAIA tests general assistant abilities {[}30{]}; PlanBench
probes planning and reasoning about change {[}31{]}; SWE-bench measures
resolution of real-world GitHub issues {[}32{]}. These benchmarks
evaluate \emph{end-to-end capability}; AgentISA's evaluation measures
\emph{subsystem behavior under controlled conditions}---call reduction,
instruction-count reduction, trace shrinkage (Section 5). The two are
complementary: instruction-level metrics can attribute end-to-end gains
to specific mechanisms when such gains are measured, and are
reproducible without live model access.

\subsection{Classical Compiler and Architecture
Techniques}\label{classical-compiler-and-architecture-techniques}

AgentISA's optimization machinery is deliberately classical:
tree-pattern instruction selection with dynamic programming, peephole
optimization, and data-flow-based dead-code elimination {[}33{]};
pipelining, hazard detection, and resource-constrained scheduling
{[}34{]}; the open, well-specified instruction-set discipline of RISC-V
{[}35{]}; type safety via progress-and-preservation {[}36{]}; formally
verified compiler correctness in the CompCert tradition {[}37{]}. Our
contribution is not a new technique but the demonstration that this
machinery transfers to LLM agent workloads, with the
semantic-preservation arguments adapted to instructions with
pre/post-conditions (Appendix C).

\subsection{LLM Serving Systems}\label{llm-serving-systems}

Below the API boundary, the serving community has made inference
dramatically cheaper: PagedAttention {[}38{]}; Orca's continuous
batching {[}39{]}; FlashAttention's IO-aware attention {[}40{]};
Mooncake's KVCache-centric disaggregation {[}41{]}; inference-cost
energy analyses {[}42{]}. This line and AgentISA are orthogonal layers
of the same stack: serving systems make each call cheaper; AgentISA
reduces how many calls an agent makes and how much redundant work each
carries---the two compound multiplicatively.

\subsection{Research Gaps and Our
Contributions}\label{research-gaps-and-our-contributions}

In summary, the literature leaves four gaps that AgentISA addresses:

\begin{enumerate}
\def\labelenumi{\arabic{enumi}.}
\item
  \textbf{No instruction-level contract for agent behavior.} Interfaces
  are standardized {[}13, 14, 18{]} and frameworks numerous {[}19, 20,
  21, 22, 23, 24{]}, but no shared instruction set spans providers and
  frameworks.
\item
  \textbf{No compiler for agent behavior.} Prompt-level compilation
  exists {[}25{]} and model-level cost reduction exists {[}28{]}, and
  \emph{caching itself is well-trodden ground}---at the response level
  (GPTCache {[}46{]}), at the KV-prefix level (RadixAttention/SGLang
  {[}47{]}), and at the provider prompt-caching level. What none of
  these layers expose is an \emph{instruction-level} contract: caching
  there is an optimization of model serving---invisible to the agent
  program's structure, and without effect on it. Instruction-level
  fusion, elimination, and scheduling---with cacheability as a declared
  property of the instruction set---have not, to our knowledge, been
  demonstrated.
\item
  \textbf{No semantics-preserving account of agent optimization.}
  Serving work verifies kernels and schedulers, not agent-level
  transformations; we provide progress-and-preservation-style arguments
  for ours (Theorems 3.2, 3.4, 3.5, 3.6).
\item
  \textbf{No reproducible instruction-level evaluation.} Existing
  benchmarks measure end-to-end capability {[}29, 30, 31, 32{]}; we
  contribute a seed-controlled evaluation: a microbenchmark suite with
  fully disclosed synthetic workloads and mock-backend instrumentation,
  extended by real-gateway and system-level validations (Section 5).
\end{enumerate}

AgentISA addresses these four gaps directly---the instruction set and
Agentic Calling abstraction (Sections 3.1--3.2), instruction-level
optimization with semantic-preservation arguments (Section 3.4, Theorems
3.2 and 3.4--3.6, Appendices B--C), trace-derived instruction learning
(Section 3.3), and a seed-controlled evaluation (Sections 5.2--5.8); the
full contribution list is in Section 1.

\section{AgentISA Design}\label{agentisa-design}

Building on Sections 1 and 2, we present the AgentISA design: the core
architecture (Section 3.1), the Agentic Calling framework (Section 3.2),
trace-based instruction learning (Section 3.3), instruction organization
and optimization (Section 3.4), and the capability-assessment layer
(Section 3.5).

\subsection{Core Architecture}\label{core-architecture}

\subsubsection{Instruction Set
Specification}\label{instruction-set-specification}

The AgentISA instruction set is formally defined as a tuple:

\textbf{Definition 3.1} (AgentISA Core): \[
\text{AgentISA} = (I, T, \Sigma, \Delta, \Phi)
\]

where:

\begin{itemize}
\tightlist
\item
  \(I\) is the set of instructions
\item
  \(T\) is the type system
\item
  \(\Sigma\) is the state space
\item
  \(\Delta\) is the transition function set
\item
  \(\Phi\) is the verification function set
\end{itemize}

Following the conventions of classical instruction set design {[}33, 34,
35{]}, we define the basic instruction format:

\begin{verbatim}
INSTRUCTION = {
    opcode: OPCODE,
    args: ARGS[],
    precondition: PRED,
    postcondition: PRED,
    context: CONTEXT
}
\end{verbatim}

The instruction set is categorized into four fundamental classes:

\begin{enumerate}
\def\labelenumi{\arabic{enumi}.}
\item
  \textbf{Control Flow Instructions}:
  \[\mathcal{C} = \{seq, branch, loop, call\}\]
\item
  \textbf{State Management Instructions}:
  \[\mathcal{S} = \{load, store, update, cache\}\]
\item
  \textbf{LLM Interaction Instructions}:
  \[\mathcal{L} = \{query, response, context\}\]
\item
  \textbf{Resource Management Instructions}:
  \[\mathcal{R} = \{allocate, free, monitor\}\]
\end{enumerate}

Two sanity results anchor the definition; neither introduces new
machinery {[}33, 36{]}: Theorem 3.1 certifies the instruction set does
not constrain what agent programs can express, and Theorem 3.2 fixes the
well-formedness obligation the offline bounded checker discharges for
pure instructions (Section 3.3.3) and the runtime brackets with
pre-dispatch context validation (Section 3.2.3).

\textbf{Theorem 3.1} (Turing Completeness): The AgentISA instruction set
is Turing complete.

\emph{Proof sketch}: a translation from WHILE programs---known Turing
equivalent---to instruction sequences, by simulation: state management
instructions encode an unbounded store, control-flow instructions encode
composition, branching, and iteration. The full construction is in
Appendix A.

\subsubsection{Type System and
Semantics}\label{type-system-and-semantics}

The type system T extends the simply-typed discipline of Pierce {[}36{]}
with dependent-type constructions (\(\Pi\), \(\Sigma\), \(\mu\)):

\textbf{Definition 3.2} (AgentISA Type System): \[
T ::= \text{Base} \mid \Pi(x:T_1).T_2 \mid \Sigma(x:T_1).T_2 \mid \mu X.T
\]

where:

\begin{itemize}
\tightlist
\item
  Base types include \{Bool, Int, String, Float\}
\item
  \(\Pi\) represents dependent function types
\item
  \(\Sigma\) represents dependent pair types
\item
  \(\mu\) represents recursive types
\end{itemize}

The semantic rules follow operational semantics:

\[
\frac{e_1 \rightarrow e_1'}{E[e_1] \rightarrow E[e_1']} \text{ (Context)} \quad
\frac{}{(\lambda x.e)v \rightarrow e[x:=v]} \text{ ($\beta$-reduction)}
\]

The pure core carries a standard type-safety argument---progress and
preservation, in the style of Pierce {[}36{]}---discharging the
well-formedness role fixed in Section 3.1 and extended to the shipped
instruction set by the bounded checker of Section 3.3.3:

\textbf{Theorem 3.2} (Type Safety for the Pure Core): Well-typed
programs in the pure core calculus of Appendix B---the
\(\lambda\)-fragment (base types, dependent \(\Pi\)/\(\Sigma\) types,
recursion) plus the composition operators applied to instructions
carrying typed pre/post-conditions---cannot ``go wrong'', under the
per-instruction validity obligation of Definition 3.5 (stated in Section
3.3.3 below; pre/post-conditions), with transition functions assumed
total on precondition-satisfying states (Appendix B.7). Effectful
instruction classes are embedded via their typed pre/post-conditions and
covered to the extent their preconditions are checkable.

\emph{Proof sketch}: By the standard two-lemma technique---progress and
preservation---as presented in Pierce {[}36{]}:

\begin{enumerate}
\def\labelenumi{\arabic{enumi}.}
\tightlist
\item
  Progress:
  \(\forall e, T.\ \vdash e : T \implies \mathrm{value}(e) \lor \exists e'.\ e \to e'\)
\item
  Preservation:
  \(\forall e, e', T.\ \vdash e : T \land e \to e' \implies \vdash e' : T\)
\end{enumerate}

Both are proved by induction over Appendix B's merged judgment system,
with a joint argument for the composition cases: progress from the
Definition 3.5 obligation plus totality, preservation from postcondition
typing, composition via the chained invariant
\(post(i_1) \Rightarrow pre(i_2)\) (Appendix B).

\subsubsection{LLM Integration
Interface}\label{llm-integration-interface}

The LLM integration interface is a protocol playing the same role for
agent backends that the Model Context Protocol plays for tool exposure
{[}18{]}: \(\mathcal{P}_{LLM} = (Q, R, C, \psi)\), where \(Q\) is the
query space (JSON Schema), \(R\) the response space (JSON Schema), \(C\)
the context-management surface (today the context-window size), and
\(\psi\) the state mapping Map\textless AgentState,
LLMContext\textgreater---corresponding to the interface conventions of
modern function-calling providers {[}3, 18{]}.

\subsection{Agentic Calling Framework}\label{agentic-calling-framework}

\subsubsection{Function Calling to Agentic Calling
Transformation}\label{function-calling-to-agentic-calling-transformation}

The transformation from function calling to Agentic Calling is
formalized as:

\textbf{Definition 3.3} (Agentic Transformation): \[
\mathcal{T}: F \rightarrow A
\]

where \(F\) is the function call space and \(A\) is the agent
instruction space.

\(\mathcal{T}\) is required to preserve call semantics; this holds by
construction for the three-step process below---the synthesized
instruction records the function signature verbatim, embeds the context
verbatim (Theorem 3.3), and delegates execution to the original
call---so we state preservation as a property of the construction rather
than a separately proved theorem:

\[
\forall f \in F: \text{sem}(\mathcal{T}(f)) \equiv \text{sem}(f)
\]

The transformation process follows three steps, mirroring the
signature-to-call conventions of function-calling interfaces {[}13,
14{]}:

\begin{enumerate}
\def\labelenumi{\arabic{enumi}.}
\tightlist
\item
  Function signature analysis
\item
  Context embedding
\item
  Instruction synthesis
\end{enumerate}

\subsubsection{Context Preservation and State
Management}\label{context-preservation-and-state-management}

State management in AgentISA follows a formal state machine model:

\[
M = (S, s_0, \delta, F)
\]

where:

\begin{itemize}
\tightlist
\item
  S: Set of states
\item
  s\(_{0}\): Initial state
\item
  \(\delta\): State transition function
\item
  F: Set of final states
\end{itemize}

Extraction-time context preservation is guaranteed by construction
(Theorem 3.3); at runtime it is only monitored:

\textbf{Theorem 3.3} (Context Preservation): Let \(\phi\) be the context
interpretation function (defined compositionally over context
\emph{contents}---the keys and values the context carries), and \(t\)
the extraction transformation that embeds a context \(c\) verbatim into
an instruction's \texttt{context} field. Then, for every \(c \in C\):

\[
\phi(t(c)) = \phi(c)
\]

\emph{Proof}: By the definition of \(t\), the extracted instruction
carries a verbatim copy of the source context; since \(\phi\) is defined
over context contents and the copy is content-identical,
\(\phi(t(c)) = \phi(c)\) holds by construction---a \emph{design
invariant of the extraction pipeline} (the implementation copies rather
than re-encodes). At runtime the shipped checks are constraint-based and
bracket the transition: \texttt{AgenticTransformer.validate\_context}
checks the registered context constraints pre-dispatch, and
\texttt{verify\_context\_preservation(before,\ after)} re-checks them
pairwise across the executed transition and reports violations
(mechanics in Section 3.2.3); neither performs structural validation of
the carried context. Extraction-time preservation is discharged by the
copy construction; preservation across an executed transition is
\emph{monitored, not guaranteed}, by the runtime checks.
\(\blacksquare\)

\subsubsection{Cross-model
Compatibility}\label{cross-model-compatibility}

Cross-model \emph{protocol} compatibility---not behavioral equivalence
of models---is achieved through an abstract interface layer
(provider-independent instruction representation with provider-specific
adapters):

\begin{verbatim}
Interface ModelAdapter {
    translate: AgentInstruction → ModelSpecificCall      // shipped: translate_instruction
    normalize: ModelSpecificResponse → AgentResponse    // shipped: normalize_response
    validateContext: Context → Boolean                  // shipped: AgenticTransformer.validate_context
}
\end{verbatim}

All three operations ship: \texttt{translate} is
\texttt{ModelAdapter.translate\_instruction} (with registered type
transformations); \texttt{normalize} is
\texttt{ModelAdapter.normalize\_response}, folding the two dominant wire
shapes---Anthropic (typed \texttt{content} blocks,
\texttt{usage.input\_tokens}/\texttt{output\_tokens}) and OpenAI
(\texttt{choices{[}0{]}.message},
\texttt{usage.prompt\_tokens}/\texttt{completion\_tokens})---into a
canonical \texttt{AgentResponse} envelope (content, tool calls, usage,
model, raw), so downstream code never branches on provider syntax; and
\texttt{validateContext} is
\texttt{AgenticTransformer.validate\_context}, the pre-execution
complement of the shipped post-hoc
\texttt{verify\_context\_preservation(before,\ after)}---it evaluates
the registered context constraints against a single state before
dispatch, applying two-argument predicates reflexively \((v, v)\)
(single-argument ones directly) and leaving constraints over absent keys
vacuously satisfied, so key \emph{deletion} goes unreported---a blind
spot analogous to non-reflexive constraints (e.g.~strict growth), which
are not detectable pre-execution.

Motivated by cross-model function-calling studies {[}17{]} and open
protocol efforts {[}18{]}---providers differ in surface syntax but can
agree on \emph{protocol shape}---the requirement is a
\textbf{protocol-level contract on the adapter layer}. For
\(\mathbb{M}\) the set of provider adapters, each \(m \in \mathbb{M}\)
provides a request normalizer \(\text{norm}_{in}(m)\) and a response
normalizer \(\text{norm}_{out}(m)\) such that:

\[
\begin{aligned}
\forall m \in \mathbb{M}: & \text{norm}_{in}(m), \text{norm}_{out}(m) \\
& \text{ are total on the supported request and response classes,}
\end{aligned}
\]

mapping every supported raw request and response into the unified
schemas, with declared types preserved and errors classified into a
common taxonomy. This is deliberately weaker than semantic-output
equivalence---\(m_1 = m_2\) would already demand determinism, which no
LLM provides---and cross-model \emph{behavioral} comparison is an
empirical question at benchmark level (Section 3.5, {[}17, 29{]}).

\subsection{Trace-Based Instruction
Learning}\label{trace-based-instruction-learning}

\subsubsection{Dialogue Pattern
Recognition}\label{dialogue-pattern-recognition}

Pattern recognition in AgentISA is frequency-based and deterministic
(action-signature grouping) rather than grammatical; it is implemented
in \texttt{agentflow.core.isa.pattern\_mining} and
\texttt{agentflow.core.isa.instruction\_learning}. To prevent
over-reading: the subsystem is a \textbf{discovery-and-replay pipeline,
not an autonomous learning loop}. Its five stages are (i) trace-derived
pattern mining (this section), (ii) dialogue-derived template extraction
(§3.3.2---templates carry no recorded instruction body and raise
\texttt{NotImplementedError} if executed), (iii) candidate validation
(support and success-rate thresholds), (iv) \emph{developer-controlled}
promotion (mining never edits the optimizer's rule set; a human decides
which candidates become instructions), and (v) composite replay (a
promoted instruction re-runs its recorded subsequence on the same
inputs; it does not generalize to unseen parameter values). Stages
(i)--(iii) and (v) ship; stage (iv) is deliberately manual.

\textbf{Definition 3.4} (Instruction Sequence Pattern): A pattern \(p\)
over the instruction alphabet \(I\) is a non-empty contiguous sequence
\(p = (i_1, \dots, i_k) \in I^k\) together with two frequency
statistics:

\begin{itemize}
\tightlist
\item
  \textbf{Support}: \[
  \begin{aligned}
  sup(p) = P(& p \text{ occurs as a contiguous} \\
  & \text{subsequence in a randomly drawn trace})
  \end{aligned}
  \]
\item
  \textbf{Confidence}:
  \(conf(p \Rightarrow j) = P(i_{t+1} = j \mid (i_{t-k+1}, \dots, i_t) = p)\)
\end{itemize}

A pattern is \emph{frequent} if \(sup(p) \geq \sigma_{min}\) and
\emph{predictive} if \(\max_j conf(p \Rightarrow j) \geq \gamma_{min}\),
for thresholds \(\sigma_{min}, \gamma_{min} \in (0, 1]\).

Pattern recognition proceeds in three stages:

\begin{enumerate}
\def\labelenumi{\arabic{enumi}.}
\item
  \textbf{Frequent-sequence mining.} Candidate subsequences are
  enumerated from execution traces with occurrence counts; successor
  statistics are accumulated alongside, and sequences below
  \(\sigma_{min}\) are pruned together with their statistics---so
  confidence is estimated only over frequent sequences. The output is
  the set of frequent, predictive patterns.
\item
  \textbf{Behavior grouping.} The shipped \texttt{PatternRecognizer}
  groups dialogues deterministically by \emph{action signature}---the
  sorted set of system-turn contents---and promotes each sufficiently
  supported bucket into a pattern. An earlier prototype (TF-IDF
  embeddings, density-based clustering, k-means) was abandoned for
  production use: non-deterministic across library versions and
  degenerate on the short traces agents actually provide (a k-means view
  over instruction-trace features remains in \texttt{pattern\_mining},
  with fixed \texttt{random\_state}). Promotions split by source:
  trace-mined promotions execute by \textbf{composite replay}---the
  learned instruction's \texttt{execute} body re-runs the recorded
  constituent sequence, with its dataflow edges, in mined order against
  a shared context, merging outputs and returning the keys its
  \texttt{produces\_X} postconditions declare---while dialogue-mined
  templates, which carry no recorded body, raise
  \texttt{NotImplementedError} rather than silently no-opping. Promotion
  is deterministic (pattern identity derives from the mined
  instruction-id sequence), and replay is plain sequential
  composition---nothing smarter---so a learned instruction reproduces
  its recorded subsequence's outputs on the same inputs but does not
  generalize to unseen parameter values; no automated
  mining-to-optimization path exists yet (promotion is a developer
  decision).
\item
  \textbf{Pattern validation.} Each candidate is checked for
  significance---support above \(\sigma_{min}\) and, for
  dialogue-derived patterns, an acceptable success rate---before
  promotion to a candidate instruction.
\end{enumerate}

This design is deliberately statistical: agent traces are short and
noisy, and frequency statistics degrade gracefully where grammar
induction would overfit. It also composes with the compiler---a
frequent, verified pattern is exactly the candidate for promotion to a
fused instruction (Section 3.4.4).

\subsubsection{Instruction Extraction
Pipeline}\label{instruction-extraction-pipeline}

The shipped extraction process instantiates the Section 3.3.1 pipeline:
frequent-sequence mining (pattern detection), deterministic
action-signature grouping (instruction formation, carrying the recorded
constituent sequence and its dataflow edges), and threshold-gated
promotion with the standard optimization passes applied to any resulting
stream (optimization and validation).

\subsubsection{Instruction Optimization and
Verification}\label{instruction-optimization-and-verification}

Optimization is coupled to verification: where verified-compilation
efforts such as {[}37{]} verify the transformations of a whole compiler,
AgentISA's verification targets a different granularity---the validity
of individual instructions and the obligations of individual rewrites:

\textbf{Definition 3.5} (Instruction Verification): An instruction i is
valid if: \[
\forall s \in S: \text{pre}(i,s) \implies \text{post}(i,\text{exec}(i,s))
\]

The shipped checker (\texttt{agentflow.core.isa.bounded\_checker})
enumerates the candidate input states derived from the declared
parameter constraints---exhaustively for finite domains such as
\texttt{bool} and enum parameters, and as a bounded, deterministically
sampled set (cap 256) otherwise---making it \emph{bounded}, not
complete. For every pure instruction it filters the state space by
well-typedness and declared preconditions, executes the instruction, and
evaluates three obligations per surviving state---progress (no raise
where the precondition holds), preservation (declared outputs present
and well-typed), and the postcondition itself. Side-effectful
instructions are never executed and are reported out of scope; their
obligation is delegated to authoring-time tests and review. Theorem 3.2
is conditional on this obligation, Theorems 3.4--3.6 carry the separate
rewrite obligations of Appendix C.7, and beyond the enumeration bound
the condition remains substantive rather than discharged (Appendices B.7
and C.7):

\[
\forall s \in S_{bounded} \subseteq S: \text{pre}(i,s) \implies \text{post}(i,\text{exec}(i,s))
\]

where \(S_{bounded}\) is the explored subset; completeness is traded for
tractability---the standard bounded-exploration compromise, with the
operational-semantics treatment of instruction meaning following
{[}36{]}.

The guarantee therefore comes in two explicitly separate layers.
\textbf{Layer 1 (theory, conditional):} Theorem 3.2 and the type-safety
results of Appendix B hold for every program whose instructions satisfy
Definition 3.5 \emph{on all of \(S\)}, under the total-definedness
premise of Appendix B.7---a mathematical statement about a proof
obligation, not about the checker. \textbf{Layer 2 (implemented
verification):} the shipped checker discharges that obligation only on
\(S_{bounded}\) (exhaustive enumeration for small finite domains;
deterministic bounded sampling with cap 256 otherwise) and only for pure
instructions; side-effectful instructions are reported out of scope.
Success of the bounded check is evidence toward, not a proof of, Layer
1's premise; no claim in this paper equates the two.

\textbf{Theorem 3.4} (Optimization Soundness): Every atomic rewrite of
the pass pipeline---fusion (Appendix C.2's rule, applied by the
combining pass), dead-code elimination, and instruction combining
(Section 3.4.4)---preserves the denotation of the instruction sequence
it transforms: \(\text{sem}(r(i)) \equiv \text{sem}(i)\), read as the
partial equivalence of C.1.

\emph{Proof sketch}: By structural induction on instruction composition.
The base case shows each atomic rewrite preserves the denotation of a
single instruction; the inductive step lifts preservation through each
composition operator of Definition 3.6. The full proof, with the
per-rewrite case analysis, is in Appendix C.

The shipped checker
(\texttt{agentflow.core.isa.bounded\_checker.verify\_instruction})
implements Definition 3.5 on the bounded state space exactly as
described above (enumerate, filter by preconditions, execute only pure
instructions); it reports violations structurally:

\begin{Shaded}
\begin{Highlighting}[]
\KeywordTok{def}\NormalTok{ verify\_instruction(instruction, pre\_condition, post\_condition):}
    \ControlFlowTok{if}\NormalTok{ instruction.side\_effects: }\ControlFlowTok{return}\NormalTok{ skipped(instruction)  }\CommentTok{\# never executed}
\NormalTok{    state\_space }\OperatorTok{=}\NormalTok{ generate\_state\_space(instruction)  }\CommentTok{\# bounded: exhaustive or seeded, cap 256}
    \ControlFlowTok{for}\NormalTok{ state }\KeywordTok{in}\NormalTok{ state\_space.states:}
        \ControlFlowTok{if}\NormalTok{ pre\_condition(state):}
\NormalTok{            next\_state }\OperatorTok{=}\NormalTok{ execute(instruction, state)}
\NormalTok{            check(progress, preservation, post\_condition, next\_state)  }\CommentTok{\# structured violation reporting}
    \ControlFlowTok{return}\NormalTok{ verdict  }\CommentTok{\# checked\_states, exhaustive, violations by kind, passed}
\end{Highlighting}
\end{Shaded}

\subsection{Instruction Organization and
Optimization}\label{instruction-organization-and-optimization}

\subsubsection{Instruction Composition}\label{instruction-composition}

The AgentISA supports hierarchical instruction composition through a
formal composition algebra:

\textbf{Definition 3.6} (Instruction Composition): Given instructions
\(i_1, i_2 \in I\), we define composition operators:

\begin{align*}
&\text{Sequential}: i_1 \circ i_2 \\
&\text{Parallel}: i_1 \parallel i_2 \\
&\text{Conditional}: i_1 \triangleright i_2 \\
&\text{Iterative}: i_1^* \\
\end{align*}

Their denotations (sequential execution; parallel interleaving under
disjoint state footprints; guarded execution of \(i_2\) on \(i_1\)'s
triggering outcome; zero-or-more iterations of \(i_1\)) are fixed in
Appendix C.1.

Composition follows algebraic properties. Properties 1 and 3 hold for
all instructions; Property 2 is a \emph{conditional} law---purity and
pairwise footprint disjointness together do \emph{not} suffice, and
idempotence of \(i_1\) must be assumed separately:

\begin{enumerate}
\def\labelenumi{\arabic{enumi}.}
\item
  \textbf{Associativity} (all instructions, up to denotation):
  \[(i_1 \circ i_2) \circ i_3 = i_1 \circ (i_2 \circ i_3)\]
\item
  \textbf{Distributivity} (conditional on idempotent \(i_1\) and
  pairwise-disjoint footprints): the law
  \[i_1 \parallel (i_2 \circ i_3) = (i_1 \parallel i_2) \circ (i_1 \parallel i_3)\]
  holds under two premises that must be stated separately. \textbf{(a)
  Idempotence of \(i_1\)}:
  \(\delta_{i_1} \circ \delta_{i_1} = \delta_{i_1}\) on the
  precondition-satisfying states---\(i_1\)'s writes must not change the
  values its own subsequent execution reads (a constant overwrite of
  declared outputs qualifies; a read-modify-write of a cell it also
  reads does not). \textbf{(b) Pairwise footprint disjointness}: the
  denotations of \(i_1, i_2, i_3\) act on pairwise disjoint state
  components (declared outputs distinct, read sets disjoint from the
  others' writes). Under (a) the second execution of \(i_1\) reproduces
  the first's writes; under (b) the other instructions cannot perturb
  \(i_1\)'s inputs between its two occurrences; disjoint updates
  commute, so both sides denote the same net transformation.

  \emph{Both premises are necessary, and neither follows from the other
  or from purity.} Purity not implying idempotence is the sharper
  failure: a pure \texttt{update(+1)} applied twice increments twice
  (purity bars \emph{external} side effects, not writes to declared
  cells). Taking \(i_1\) as such an update and \(i_2 = i_3 = \epsilon\)
  (the identity, empty footprint, so (b) holds vacuously), the two sides
  execute \(i_1\) once and twice---from \(x = 0\) they yield \(x = 1\)
  and \(x = 2\); only premise (a) rules this out. Conversely, without
  (b), a write of \(i_2\) into a cell \(i_1\) reads changes \(i_1\)'s
  second execution even when \(i_1\) is idempotent.
\item
  \textbf{Identity}: \[i \circ \epsilon = \epsilon \circ i = i\]

  where \(\epsilon\) is the empty (identity) instruction.
\end{enumerate}

\textbf{Remark} (implementation consequence): distributivity also fails
for side-effectful instructions, for the same double-execution reason.
Because the law is doubly conditional, the shipped compiler does not
rely on it at all: parallel grouping is gated on purity plus disjoint
dependencies and disjoint outputs, \texttt{side\_effects=True}
instructions are excluded from all such rewrites, and no
distributivity-based simplification is applied. The law is stated for
the algebra's completeness, not as a license for rewrites.

\subsubsection{Instruction Pipeline
Design}\label{instruction-pipeline-design}

We formalize the instruction pipeline as a directed acyclic graph (DAG):

\textbf{Definition 3.7} (Pipeline DAG): \[
P = (V, E, \lambda, \delta)
\]

where:

\begin{itemize}
\tightlist
\item
  V: Pipeline stages
\item
  E: Dependencies between stages
\item
  \(\lambda\): Stage latency function
\item
  \(\delta\): Data transfer function
\end{itemize}

The pipeline's stages (FETCH \(\rightarrow\) DECODE \(\rightarrow\)
SCHEDULE \(\rightarrow\) EXECUTE \(\rightarrow\) COMMIT, with hazard
detection, dependency tracking, and resource allocation maps; listing in
Appendix D.11 of the supplementary material) are organized around this
DAG.

\textbf{Theorem 3.5} (Pipeline Correctness): For any valid instruction
sequence \(s\) and any schedule \(S\) satisfying conditions S1-S4 of
Appendix C.5---each instruction dispatched exactly once, dependencies
respected, dispatch groups pairwise non-interfering, commit in program
order---pipeline execution preserves program semantics.

\emph{Proof}: By induction over the dispatch groups of \(S\) in commit
order, under the invariant that the architectural state after committing
each prefix of groups equals the sequential state after the same
instructions; S1--S4 supply exactly-once dispatch, dependency order,
group commutativity, and in-order commit (details in Appendix C.5).

\subsubsection{Instruction Selection and
Scheduling}\label{instruction-selection-and-scheduling}

Instruction selection follows a tree-pattern matching approach:

\textbf{Definition 3.8} (Selection Pattern): \[
\pi = (\tau, C, \text{cost}, \text{code})
\]

where:

\begin{itemize}
\tightlist
\item
  \(\tau\): a tree pattern over the IR (unrelated to the type-system
  symbol \(T\) of Definition 3.2)
\item
  \(C\): context constraints on the surrounding state
\item
  cost: Cost function
\item
  code: Generated code
\end{itemize}

Given an IR tree \(t\) to cover, the selection process is formalized as:

\[
\text{Select}(t) = \argmin_{\pi \in \Pi} \{\text{cost}(\pi) \mid \text{matches}(\pi.\tau, t) \land \text{satisfies}(t, \pi.C)\}
\]

Scheduling is formalized as the classical delay-minimization problem
under dependency and capacity constraints (formulation in Appendix E.1
of the supplementary material): minimize the total stall
\(\sum_i \text{delay}(i)\) subject to
\(\text{start}(i_2) \geq \text{end}(i_1)\) for every dependency pair and
\(\sum_{i \in \text{active}(t)} \text{res}(i) \leq \text{Capacity}\) at
every time \(t\)---the specified, not shipped, scheduler's objective.

\subsubsection{Static and Dynamic
Optimization}\label{static-and-dynamic-optimization}

\textbf{Static optimization.} Static optimization employs several
techniques:

\begin{enumerate}
\def\labelenumi{\arabic{enumi}.}
\item
  \textbf{Peephole Optimization}: \[
  \text{PH}(i_1,...,i_n) = \begin{cases}
  \text{opt}(i_1,...,i_n) & \text{if pattern exists} \\
  i_1,...,i_n & \text{otherwise}
  \end{cases}
  \]
\item
  \textbf{Dead Code Elimination}: \[
  \text{Live}(i) = \text{Use}(i) \cup (\text{Live}_{\text{out}}(i) - \text{Def}(i))
  \] The two shipped passes approximate this fixpoint at different
  conservatism levels: the \emph{compiler} pass (formal streams) removes
  a side-effect-free instruction when no live consumer depends on its id
  or any declared output, while the \emph{trace optimizer} pass
  (executable instructions) keeps any instruction declaring at least one
  output---executable traces lack the formal dependency annotations that
  make output-consumption analysis exact. Section 5.3 exercises the
  former; Section 5.4 the latter.
\item
  \textbf{Instruction Combining} (the named candidate rules below are
  specified and not registered---Section 4.1.3; the shipped optimizer's
  built-in pairwise compatibility combining is active, and it is this
  built-in combining that the hot-path trace re-optimization of Section
  5.4 applies):

\begin{verbatim}
Combine Rules {
  // LOAD/STORE/MOVE are textbook mnemonics, illustrative
  // only; they are not shipped opcode names
  LOAD followed by STORE → MOVE
  Multiple CONTEXT_UPDATE → Single UPDATE
  Sequential QUERY → Batch QUERY
}
\end{verbatim}
\end{enumerate}

\textbf{Dynamic optimization.} Dynamic optimization in the shipped
runtime is \textbf{repetition-triggered and threshold-based}; it does
not (yet) respond to load or performance signals:

\begin{enumerate}
\def\labelenumi{\arabic{enumi}.}
\item
  \textbf{Hot Path Detection} (implemented): the
  \texttt{DynamicOptimizer} keeps a monotone counter per exact
  instruction sequence (\texttt{path\_id} = the ordered join of
  instruction ids); a path is flagged hot when its count exceeds a
  configurable threshold \emph{and} its most recent recorded success
  rate is at least a stability bound (0.9 by default). There is no
  windowing and no load signal: \[
  \text{Hot}(p) :\!\!\iff \text{count}(p) > \theta \land \text{success}(p) \geq \alpha
  \]
\item
  \textbf{Trace Optimization} (implemented): once hot, the cached
  sequence is re-optimized by the static pass machinery---peephole,
  dead-code elimination, instruction combining---bounded to a fixed
  number of passes; the runtime then serves the shrunken form for
  subsequent executions (demonstrated end-to-end in Section 5.8).
\item
  \textbf{Trace Formation and Adaptive Recompilation} (specified, not
  implemented): sub-trace extraction from observed frequency profiles,
  and performance-triggered recompilation (\texttt{should\_recompile}:
  execution count above threshold \emph{and} observed performance below
  expectation) remain design-level. The shipped runtime re-optimizes the
  whole cached sequence once per recording after the threshold is
  crossed, with no performance feedback loop.
\end{enumerate}

Two known consequences of the shipped design: because \texttt{path\_id}
requires exact sequence identity, trace diversity
(similar-but-not-identical sequences) never accumulates counts and the
mechanism stays cold; and hot-path state (sequences, counters,
statistics) is unbounded. Both are discussed with their remedies in
Section 6.1.

\subsubsection{Verification and
Analysis}\label{verification-and-analysis}

Optimization correctness is argued through:

\textbf{Theorem 3.6} (Pass-Pipeline Soundness, extending Theorem 3.4):
For any optimization transform T composed of the pass pipeline's
rewrites, \[
\forall p \in \text{Program}: \text{sem}(T(p)) \equiv \text{sem}(p)
\]

\textbf{Proof Sketch}:

\begin{enumerate}
\def\labelenumi{\arabic{enumi}.}
\tightlist
\item
  Show semantic preservation for each atomic transformation (the
  atomic-rewrite lemma, Appendix C.2)
\item
  Prove composition of transformations preserves semantics (the
  compositionality lemma, Appendix C.3)
\item
  Note that hazard-free pipeline execution is \emph{not} part of this
  theorem---it is Theorem 3.5 (Appendix C.5), invoked separately; this
  theorem's contribution is the rewrite layer (steps 1--2).
\end{enumerate}

\textbf{Performance Model} (specified; the benchmarks of Section 5
measure latency and token counts directly and do not evaluate it): \[
\text{Perf}(p) = \sum_{i \in p} (\text{base\_cost}(i) + \text{overhead}(i) - \text{optimization\_gain}(i))
\]

where optimization\_gain is calculated as: \[
\text{optimization\_gain}(i) = \sum_{o \in \text{Opts}} \text{benefit}(o,i) \times \text{application\_rate}(o,i)
\]

\subsection{Capability Assessment}\label{capability-assessment}

Agents are heterogeneous: an agent may be strong at reasoning but weak
at tool use, or vice versa. To make such differences explicit---and to
give the multi-agent coordinator of Section 5.6 a principled basis for
agent selection---AgentISA includes a capability assessment subsystem,
implemented as \texttt{CapabilityProfile} in
\texttt{agentflow.core.isa.capability}.

\textbf{Definition 3.9} (Capability Profile): A capability profile of an
agent \(a\) maps each capability domain \(D\) (the shipped enum carries
eight domains: reasoning, memory, generation, optimization,
specialization, coordination, security, robustness) to a score--weight
pair: \[
\text{Cap}(a) = \{(\text{dom}_k, s_k, w_k) \mid \text{dom}_k \in D\}
\]

where each domain carries a score \(s_k \in [0, 1]\) and a weight
\(w_k \geq 0\) reflecting the domain's importance in the current
deployment context.

\textbf{Implementation status.} The aggregation and selection mechanisms
ship (\texttt{CapabilityProfile}, \texttt{get\_overall\_score}, and the
coordinator's threshold selection) and are functionally tested in
Section 5.6; the instruction-level metric extraction is a specified
stub---the seventeen \texttt{\_measure\_*} functions return 0.0---so
\texttt{assess\_capabilities} has no production call site yet, and
Section 5.6's tests inject \texttt{DomainScore} values directly. The
pipeline's \emph{mechanics} are pinned by unit tests
(\texttt{tests/core/isa/test\_capability\_pipeline.py}): end-to-end
assessment over a real instruction list produces well-formed results for
all eight domains with stub-zero scores; the domain aggregation is the
confidence-weighted mean with the weight deliberately absent (guarding
against a double-weighting defect observed during development---an
earlier form folded the domain weight into \(s_k\) as well, capping an
assessment below the coordinator's 0.7 threshold); the top-level
weighted mean applies weights exactly once.

Assessment proceeds in three steps:

\begin{enumerate}
\def\labelenumi{\arabic{enumi}.}
\item
  \textbf{Domain metric computation.} For each domain,
  \texttt{assess\_capabilities} is \emph{specified} to compute a set of
  \texttt{CapabilityMetric} values from the agent's instruction
  set---e.g., for the reasoning domain, metrics derived from the
  prevalence and complexity of reasoning-type instructions. As stated
  above, this extraction is a specified stub in the shipped code, so
  exercised assessments derive domain scores from injected metrics.
\item
  \textbf{Domain aggregation.} Metrics are aggregated into a per-domain
  \texttt{DomainScore} as a confidence-weighted mean, matching the
  shipped \texttt{\_compute\_weighted\_score}: \[
  s_k = \frac{\sum_{m \in M_k} v_m \cdot \text{conf}_m}{\sum_{m \in M_k} \text{conf}_m}
  \]
\end{enumerate}

where \(M_k\) is the metric set of domain \(k\), \(v_m\) the metric
value, and \(\text{conf}_m\) its confidence; an empty metric set scores
\(0\), and zero total confidence degenerates to the unweighted mean. The
domain weight \(w_k\) is deliberately absent here---it enters only once,
at the top level (an earlier shipped form folded it into \(s_k\) as
well, double-counting it; Section 5.6's tests were unaffected because
they inject \texttt{DomainScore} values directly).

\begin{enumerate}
\def\labelenumi{\arabic{enumi}.}
\setcounter{enumi}{2}
\tightlist
\item
  \textbf{Overall score.} \texttt{get\_overall\_score} computes the
  weighted mean across assessed domains: \[
  S(a) = \frac{\sum_{k} w_k \, s_k}{\sum_{k} w_k}
  \]
\end{enumerate}

with \(S(a)\) falling back to the unweighted mean when all weights are
zero.

Three properties of this design deserve emphasis. First, it is
\textbf{static rather than live-probing}: assessment is designed to be
cheap, deterministic, and side-effect-free, though the shipped metric
extraction is the stub disclosed above and exercised scores are
injected. Second, it is \textbf{compositional}: because scores summarize
instructions, fusing or eliminating instructions during compilation
(Section 3.4) changes the profile predictably, and a production
deployment would re-assess after optimization (the shipped coordinator
does not yet do so). Third, it directly serves coordination: the
agent-selection predicate in Section 5.6 chooses the highest-scoring
agent above a capability threshold for each task, subject to the
security policies of the \texttt{SecurityMonitor}---with the tests
injecting scores directly, as disclosed above.

We state the scope of this subsystem honestly:
\texttt{CapabilityProfile} assesses \emph{declared instruction-level
capability}, which correlates with---but does not measure---end-to-end
task performance. Cross-model behavioral comparison at benchmark level
is the domain of evaluation suites such as BFCL {[}17{]} and AgentBench
{[}29{]}; the two views are complementary.

\section{Implementation}\label{implementation}

This section describes the implementation status of Section 3's
components, with emphasis on the compiler infrastructure and
optimization techniques of Section 3.4.

\subsection{Compiler Infrastructure}\label{compiler-infrastructure}

\subsubsection{Instruction Selection
Implementation}\label{instruction-selection-implementation}

We \emph{specify} the tree-pattern matching approach of Section 3.4.3
(Definition 3.8) as dynamic programming over the IR tree, in the
classical manner of bottom-up instruction selection {[}33{]}.
\textbf{The shipped compiler does not yet implement this selector}: its
current selection step is a heuristic matcher---name-overlap scoring
over the instruction catalogue---whose compact listing is deferred to
the source (\texttt{agentflow.core.isa.instruction\_selector}); the DP
listing below is the \emph{specified} optimal selector, and none of the
benchmark numbers in Section 5 depend on tree-pattern selection. The
specified algorithm---the MATCH-TREE dynamic program over the IR
tree---is listed in Appendix D.1.

\textbf{Theorem 4.1} (Selection Optimality, specified algorithm): Given
a set of context-local patterns P and an IR tree T, the MATCH-TREE
algorithm computes the minimal covering cost; the covering itself is
read off the \texttt{selected} table afterwards, by walking from the
root and skipping subtrees spanned by the chosen pattern.

\emph{Proof sketch}: by induction on tree height; the recurrence sums
only over children not covered by the candidate pattern, so minima
compose (full argument, with the context-locality side condition, in
Appendix D.1). This theorem is a property of the specified algorithm,
not of the shipped heuristic selector; no optimality claim is made for
the implementation.

\subsubsection{Pipeline Scheduling
Implementation}\label{pipeline-scheduling-implementation}

We \emph{specify} the pipeline scheduler following the
instruction-scheduling formulation of Section 3.4.3 (Deps and resource
capacity) and the correctness conditions of Theorem 3.5 (Section 3.4.2),
with hazard detection and resource-constrained issue in the style of
classical pipeline scheduling {[}34{]}. \textbf{This cycle-accurate
scheduler is design-level}: the executed runtime path is the sequential
and parallel-grouped dispatch produced by the compiler's grouping pass
(Section 4.1.3); no benchmark in Section 5 exercises a cycle-accurate
schedule. The specified scheduler is listed in Appendix D.2.

Schedule verification checks the dependency-ordering condition (S2) of
the formal model (Theorem 3.5, Appendix C.5), with the edge convention
of Section 3.4.3---\((i_1, i_2) \in \text{Deps}\) means \(i_1\) must
finish before \(i_2\) starts. The check as written,

\[
\forall (i_1, i_2) \in \text{Deps}: \text{slot}(i_2) > \text{slot}(i_1),
\]

is equivalent to \(start(i_2) \geq end(i_1)\) only under the
\textbf{unit-slot model} the shipped scheduler assumes: one slot per
cycle, every instruction occupying exactly one slot, so
\(start(i) = \text{slot}(i) \cdot c\) and
\(end(i) = (\text{slot}(i) + 1) \cdot c\) for cycle length \(c\). If
instructions with multi-cycle latencies were admitted, the slot
comparison would have to be replaced by an explicit \(start\)/\(end\)
comparison (or slots generalized to intervals); the pipeline scheduler
that would need this remains the specified component of Section 4.1.2.
S1 (exactly-once dispatch) and S4 (in-order commit) hold by construction
of the scheduling loop, and S3 (group non-interference) is enforced by
the grouping check.

\subsubsection{Static Optimization
Implementation}\label{static-optimization-implementation}

We implement the optimization passes defined in Section 3.4.4. The
passes that fire in the benchmarks are fusion (the instruction-combining
pass of Section 3.4.4, listed in Appendix D.10 of the supplementary
material), dead-code elimination (the shipped two-level design of
Section 3.4.4), and parallel grouping (described with its measurement in
Section 5.3). The shipped rule set registers no peephole patterns, so
every peephole obligation below is currently dormant but stated for
completeness:

\begin{enumerate}
\def\labelenumi{\arabic{enumi}.}
\tightlist
\item
  \textbf{Peephole Optimization} (shipped class with an empty pattern
  table; listing in Appendix D.8):
\end{enumerate}

Each peephole rule carries a correctness \emph{obligation}, discharged
as part of rule definition rather than by the running compiler:

\begin{align*}
&\forall p \in \text{Patterns}, w \in \text{Instructions}^{\text{len}(p)}: \\
&\text{matches}(p, w) \implies \text{sem}(\text{replace}(p, w)) \equiv \text{sem}(w)
\end{align*}

As stated above, the shipped rule set is empty, so the equivalence
checker is dormant; Section 4.3.1 specifies the symbolic-execution
checker a populated rule set would require.

\begin{enumerate}
\def\labelenumi{\arabic{enumi}.}
\setcounter{enumi}{1}
\tightlist
\item
  \textbf{Instruction Combining}: the shipped pairing pass groups
  adjacent compatible instructions and verifies each merged group; its
  full listing is in Appendix D.10 of the supplementary material, and in
  the source module
  \texttt{agentflow.core.isa.instruction\_optimization}.
\end{enumerate}

\subsection{Runtime System}\label{runtime-system}

\subsubsection{Dynamic Pipeline Control}\label{dynamic-pipeline-control}

Design of the dynamic pipeline controller (\textbf{specified, not
implemented}---the shipped dynamic machinery is the counter/threshold
hot-path detector and the trace re-optimization described in Section
3.4.4 (both discussed in Section 4.2.2; measured in Sections 5.4 and
5.8); a skeleton module exists but is not wired into the runtime; the
listing is in Appendix D.3):

Pipeline performance is modeled with the classical pipeline
bound---throughput is the reciprocal of the slowest stage's service
time:

\[
\text{Throughput} = \frac{1}{\max_i\, \text{service\_time}(\text{stage}_i)}
\]

(a model of the specified controller; nothing in Section 5 measures it)

\subsubsection{Dynamic Optimization
Implementation}\label{dynamic-optimization-implementation}

The \emph{shipped} mechanism is the counter/threshold detector described
in Section 3.4.4: a monotone count per exact \texttt{path\_id}, crossed
against a configurable threshold with a last-recorded-success-rate gate.
What the benchmarks measure (Sections 5.4 and 5.8) is this mechanism.
The following is the \emph{specified} detector---exponential moving
average over a window, frequency normalized by recent execution share,
and trace formation from hot-path regions; it is \textbf{not yet
implemented}. Together with the block-granular trace formation below it
addresses workloads with trace diversity; the detector itself remains
exact-path-keyed (the limitation of Section 3.4.4---that similar but
non-identical sequences never accumulate a count---still applies to it):

\begin{enumerate}
\def\labelenumi{\arabic{enumi}.}
\tightlist
\item
  \textbf{Hot Path Detection} (specified; listing in Appendix D.4):
\end{enumerate}

The decision rule of the pseudocode is the threshold predicate
IS-PATH-HOT (frequency above \(\theta\) and stability above the
stability threshold); it computes no score. A \emph{ranking
variant}---specified, not part of the decision rule---can order
candidate paths by a heuristic score blending the already-normalized
frequency with recent stability through a static weight \(\beta\)
(formula in Appendix E.1); it is a heuristic score, not a conditional
probability.

\begin{enumerate}
\def\labelenumi{\arabic{enumi}.}
\setcounter{enumi}{1}
\tightlist
\item
  \textbf{Trace Formation} (specified; listing in Appendix D.5):
\end{enumerate}

Trace validity is \emph{defined} by the conjunction of three properties
(H the loop header of the region, tail the last block of the trace);
VERIFY-TRACE-PROPERTIES checks them:

\[
\begin{aligned}
\text{Valid}(\text{trace}) \iff {}& \text{Entry}(\text{trace}) \text{ dominated by } H \\
\wedge\;& \text{Exit}(\text{trace}) \text{ post-dominates tail} \\
\wedge\;& \text{NoSideExit}(\text{trace}) \text{ except tail}
\end{aligned}
\]

\subsection{Optimization Verification}\label{optimization-verification}

\subsubsection{Static Analysis}\label{static-analysis}

We \emph{specify} verification of the static optimizations from Section
3.4.4 (a symbolic-execution equivalence checker over transformation path
pairs). As stated in Section 4.1.3, the shipped peephole table is empty
and this obligation is dormant; the soundness theorems (Theorems
3.4--3.6) are conditional on discharging it when rules are introduced
(Appendix C.7). The specified checker is listed in Appendix D.6.

The verification follows the formal model:

\[
\text{Equiv}(p_1, p_2) \iff \forall \sigma \in \Sigma: \llbracket p_1 \rrbracket \sigma = \llbracket p_2 \rrbracket \sigma
\]

\subsubsection{Dynamic Analysis}\label{dynamic-analysis}

The runtime \emph{monitors} optimizations through a single mechanism (no
result-correctness oracle is defined at runtime---correctness of
LLM-instruction results is not decidable in general and is delegated to
the authoring-time verification of Section 3.3.3; the monitor below is a
specified, not shipped, component):

\begin{enumerate}
\def\labelenumi{\arabic{enumi}.}
\tightlist
\item
  \textbf{Performance Monitoring} (specified; listing in Appendix D.7):
\end{enumerate}

The evaluation protocol of Section 5.1.3 reports per-metric mean \(\pm\)
standard deviation and makes no inferential (hypothesis-test) claims; no
aggregate weighted-mean score is computed.

\subsubsection{Empirical Results}\label{empirical-results}

The implemented optimizations are validated by the microbenchmark suite
of Section 5, which drives the real compiler, cache, and
optimization-manager code paths against a mock LLM backend. We keep all
numbers in Section 5 and state only the mechanism-to-experiment map
here:

\begin{table}[t]
\caption{Mechanism-to-experiment map; every measurement is a
controlled-setting result, detailed in Section 5.}
\small\setlength{\tabcolsep}{3pt}
\begin{tabular}{@{}
  >{\raggedright\arraybackslash}p{(\linewidth - 2\tabcolsep) * \real{0.5149}}
  >{\raggedright\arraybackslash}p{(\linewidth - 2\tabcolsep) * \real{0.4851}}@{}}
\toprule\noalign{}
\begin{minipage}[b]{\linewidth}\raggedright
Mechanism (implementation status)
\end{minipage} & \begin{minipage}[b]{\linewidth}\raggedright
Experiment
\end{minipage} \\
\midrule\noalign{}
Result caching (implemented) & Mock-backend caching, 3 skew levels
(Section 5.2) \\
Fusion + DCE static passes (implemented, fire in benchmarks) & Static
compilation, 4 levels, 3 lengths (Section 5.3) \\
Hot-path detection and adoption (implemented) & Dynamic trace
re-optimization (Section 5.4) \\
Pass-level attribution (paired + independent) & Ablation (Section
5.5) \\
Coordination + security & 43 functional tests (Section 5.6) \\
Cache, compilation, hot-path adoption (live gateway) & Real-gateway
validations (Sections 5.7, 5.7.1) \\
Whole-system composition & Four-task characterization (Section 5.8) \\
\bottomrule\noalign{}
\end{tabular}
\end{table}

All measurements are controlled-setting results on synthetic workloads
(simulated 80 ms per call); production end-to-end claims are out of
scope, and Section 5.9 states the threats to validity.

\section{Empirical Evaluation}\label{empirical-evaluation}

We evaluate AgentISA with a reproducible microbenchmark suite
(\texttt{examples/benchmarks/isa\_microbench.py}) and a multi-agent
functional test suite, under the scope fixed at the end of Section 4.3:
Sections 5.2--5.5 run against the mock backend; Sections 5.7--5.8 extend
validation to a live LLM gateway and four system-level tasks; Section
5.9 discusses threats to validity.

\subsection{Methodology}\label{methodology}

\subsubsection{Instrumentation: Mock LLM
Backend}\label{instrumentation-mock-llm-backend}

All experiments route model inference through a \texttt{MockLLMBackend}
that counts each \texttt{infer} call as one simulated LLM API round trip
and charges a fixed 80 ms simulated latency per call to a simulated
clock. It never sleeps and never contacts a live model, so every ``API
call'' is observable, deterministic given a seed, and reproducible
without network access or inference spend. Everything \emph{around} the
backend is real: the cache, compiler, and optimizer are the production
\texttt{agentflow.core.isa} implementations (\texttt{InstructionCache},
\texttt{InstructionCompiler}, \texttt{OptimizationManager}), exercised
through the same code paths a deployment uses.

\subsubsection{Workloads}\label{workloads}

Workloads are synthetic and seed-controlled:

\begin{itemize}
\tightlist
\item
  \textbf{Caching}: 2,000 tasks from a discrete distribution over 50
  instruction templates \(\times\) 20 parameter values, popularity
  Zipf-distributed with exponent \(s \in \{0.5, 1.0, 1.5\}\) (larger
  \(s\) concentrates popularity on fewer tasks, increasing reuse
  potential). Three arms share the same instrumentation (Section 5.2):
  no cache, the \texttt{InstructionCache}, and a naive app-level
  exact-response cache (a plain dict keyed on the request payload).
\item
  \textbf{Static compilation}: random sequences of length 50, 200, and
  1,000 from \texttt{synthetic\_formal\_sequence}
  (\texttt{examples/benchmarks/isa\_microbench.py}; short enough to
  audit line by line). Each of the first \(n-1\) slots is drawn
  independently: with probability 0.25 a \emph{dead} instruction
  (side-effect-free, \textbf{no declared outputs, hence no consumers
  anywhere in the sequence}---removable by construction, planted not
  inferred); with probability 0.55 a \emph{pure} instruction
  (side-effect-free, one declared output \(v_i\), consumed \textbf{only
  by the final sink}, fusable next to a same-type pure neighbor);
  otherwise \emph{side-effectful} (always live, never fused or removed).
  The \(n\)-th slot is a side-effectful \emph{final sink} depending on
  exactly the set of all pure outputs, so every pure instruction has a
  downstream consumer and no pure output is deleted as dead. The
  attributes are disjoint and jointly exhaustive per slot; the passes
  find the planted DCE/fusion targets, so the measured reductions are
  properties of this workload family, not of arbitrary agent programs.
\item
  \textbf{Dynamic optimization}: instruction traces of length 100
  recorded for up to 200 iterations per seed, hot-path detection active
  (threshold 50 recordings); the loop terminates at the
  trigger---iteration 50 in all seeds---so each seed contributes exactly
  51 recordings (the 50 counted plus the trigger).
\item
  \textbf{Ablations}: two arms. The primary arm is \emph{paired}: it
  re-compiles the static-compilation sequences (identical seed base
  2,000) with fusion and dead-code-elimination toggled individually and
  jointly---a same-input decomposition of the 34.2\% headline. The
  secondary arm is an \emph{independent} draw, same distribution but
  seed base 4,000, a fresh-sequence replication (all-passes total
  35.0\%, within one standard deviation of the headline).
\end{itemize}

\subsubsection{Protocol and Statistics}\label{protocol-and-statistics}

Each synthetic-workload configuration (Sections 5.2--5.5) is run with 5
random seeds; we report mean \(\pm\) sample standard deviation. These
are \textbf{descriptive statistics only}: we make no hypothesis-test
claims (no p-values, no significance tests) anywhere in this paper. Five
seeds characterize variability across the tested workload instances, not
statistical significance or population-level generality. Section 5.2's
naive response-cache baseline shares the mock-backend instrumentation; a
formal \emph{inferential} comparison against full external systems
remains out of scope, and no such claim is made.

\subsubsection{Metric Terminology}\label{metric-terminology}

Several distinct quantities must not be conflated; each is defined once
here:

\begin{itemize}
\tightlist
\item
  \textbf{Effective-instruction reduction}---fewer \emph{effective}
  instructions in the compiled program (Sections 5.3/5.5); an
  instruction-stream property, not execution time.
\item
  \textbf{Dispatch-count reduction}---fewer dispatch operations at
  runtime, including parallel grouping's folding (Section 5.3); a
  runtime count.
\item
  \textbf{Trace shrinkage}---shorter traces after hot-path
  re-optimization (Section 5.4); the optimized trace's executed-state
  identity is verified separately there.
\item
  \textbf{Cache hit rate / call reduction}---fraction of calls served
  from the result cache; here the call-reduction ratio equals the hit
  rate by construction (Section 5.2).
\item
  \textbf{Simulated latency saving}---avoided-call counts \(\times\) the
  mock's fixed 80 ms (Section 5.2); derived, not measured.
\item
  \textbf{Measured latency/token saving}---wall-clock time and token
  counts against the real gateway (Sections 5.7--5.7.1); measured.
\end{itemize}

For context preservation (Theorem 3.3): \emph{extraction-time context
preservation} is the copy invariance of the extraction step;
\emph{runtime context validation} is the monitor's pre/post-dispatch
constraint checking---different properties, different coverage (Section
3.2.2 discloses the runtime side's deletion blind spot).

\subsection{Experiment 1: Result
Caching}\label{experiment-1-result-caching}

\textbf{Research question.} How many simulated LLM API calls does
AgentISA's result caching eliminate, as a function of workload skew?

We execute 2,000 tasks per seed, with caching disabled and with the
\texttt{InstructionCache} enabled, measuring simulated API calls at the
backend; the reduction ratio and hit rate coincide by construction
(every avoided call is exactly one cache hit).

\textbf{Cache-key contract.} The shipped result cache keys an entry on
the (instruction, params) pair---deliberately \emph{not} on model
identity/version, prompt-template version, tool schema, decoding
parameters, or ambient context. Soundness requires those factors held
fixed, which every benchmark here guarantees by construction (single
alias, fixed templates and decoding, fresh cache per arm); Section
5.7.1's classifier-keyed step hashes exactly the input dictionary its
prompt is built from, under the same discipline. A deployment changing
any of these factors must invalidate or namespace the cache (the
production cache isolates by TTL, LRU capacity, and never-cached
failures but does not detect semantic-configuration drift); cache
validity across \emph{varying} model versions or prompts is future work,
not exercised by any number in this paper.

\begin{table}[t]
\caption{Simulated API calls with and without the result cache (2,000
tasks per seed, 5 seeds, mean \(\pm\) sample standard deviation;
reduction equals hit rate by construction; the third, naive
response-cache arm is reported in the text below).}
\small\setlength{\tabcolsep}{3pt}
\begin{tabular}{@{}
  >{\raggedright\arraybackslash}p{(\linewidth - 6\tabcolsep) * \real{0.2537}}
  >{\raggedright\arraybackslash}p{(\linewidth - 6\tabcolsep) * \real{0.2388}}
  >{\raggedright\arraybackslash}p{(\linewidth - 6\tabcolsep) * \real{0.2090}}
  >{\raggedright\arraybackslash}p{(\linewidth - 6\tabcolsep) * \real{0.2985}}@{}}
\toprule\noalign{}
\begin{minipage}[b]{\linewidth}\raggedright
Zipf exponent \(s\)
\end{minipage} & \begin{minipage}[b]{\linewidth}\raggedright
Calls (no cache)
\end{minipage} & \begin{minipage}[b]{\linewidth}\raggedright
Calls (cached)
\end{minipage} & \begin{minipage}[b]{\linewidth}\raggedright
Reduction = hit rate
\end{minipage} \\
\midrule\noalign{}
0.5 & 2,000 & 810.4 \(\pm\) 5.7 & 59.5\% \(\pm\) 0.3\% \\
1.0 & 2,000 & 642.0 \(\pm\) 3.8 & 67.9\% \(\pm\) 0.2\% \\
1.5 & 2,000 & 408.0 \(\pm\) 9.4 & 79.6\% \(\pm\) 0.5\% \\
\bottomrule\noalign{}
\end{tabular}
\end{table}

\textbf{Naive response-cache baseline.} A third arm runs an app-level
\emph{response} cache keyed on the exact request payload string---what
an exact-match cache in front of the backend (GPTCache's exact mode, or
an HTTP response cache) keys on---implemented as a plain dict with no
AgentISA machinery. Across the three skews it issues exactly the same
call counts as the table's cached arm: \textbf{parity with the
\texttt{InstructionCache} to the digit}, because the (instruction,
params) key is already canonical on this workload. Parity is the
finding, stated plainly: at the level of call reduction the
instruction-level cache adds nothing over exact-response caching; what
it adds is the compile-time transformations of Sections 5.3 and
5.5---fusion and dead-code elimination, reducing the program before any
request is issued, which a response cache cannot express because it sees
request strings, not instruction semantics. Semantic/embedding caching
(GPTCache's approximate mode {[}46{]}) remains out of scope: no
head-to-head against approximate matching is claimed.

\textbf{Interpretation.} The reduction grows monotonically with skew,
59.5\% at \(s = 0.5\) to 79.6\% at \(s = 1.5\): the savings' upper bound
is the probability mass of repeated tasks, which the Zipf exponent
directly controls. In simulated latency the savings are the avoided-call
counts \(\times\) 80 ms (\(s = 0.5/1.0/1.5\): 95,168 \(\pm\) 455 /
108,640 \(\pm\) 305 / 127,360 \(\pm\) 750 ms)---derived, not separately
measured.

Two caveats keep this honest. First, these are \emph{simulated} calls
against a mock backend: the reduction is exact for this workload family,
but real deployments face cache-invalidation, prompt-context, and
embedding-variability effects the mock abstracts away. Second, the
absolute number depends on reuse structure; workloads without repetition
would see zero benefit. The experiment establishes the \emph{mechanism
and its skew sensitivity}, not a universal constant.

\subsection{Experiment 2: Static
Compilation}\label{experiment-2-static-compilation}

\textbf{Research question.} How much does the compiler reduce
instruction sequences, and what do the different optimization levels
contribute?

We compile random sequences at four optimization levels and measure two
quantities, kept explicit because they are easy to conflate:
\textbf{effective instruction count}---instructions remaining after
fusion and dead-code elimination, a program-size measure---and
\textbf{dispatch count}---independent dispatch operations the runtime
issues, where parallel grouping can fold several independent
instructions into one dispatch \emph{without changing the effective
instruction count}.

\begin{table}[t]
\caption{Effective-instruction-count reduction by optimization level and
sequence length (5 seeds, mean \(\pm\) sample standard deviation;
planted dead 0.25 / pure 0.55 workload).}
\small\setlength{\tabcolsep}{3pt}
\begin{tabular}{@{}
  >{\raggedright\arraybackslash}p{(\linewidth - 6\tabcolsep) * \real{0.2143}}
  >{\raggedright\arraybackslash}p{(\linewidth - 6\tabcolsep) * \real{0.1571}}
  >{\raggedright\arraybackslash}p{(\linewidth - 6\tabcolsep) * \real{0.4286}}
  >{\raggedright\arraybackslash}p{(\linewidth - 6\tabcolsep) * \real{0.2000}}@{}}
\toprule\noalign{}
\begin{minipage}[b]{\linewidth}\raggedright
Sequence length
\end{minipage} & \begin{minipage}[b]{\linewidth}\raggedright
Level
\end{minipage} & \begin{minipage}[b]{\linewidth}\raggedright
Effective-instruction reduction
\end{minipage} & \begin{minipage}[b]{\linewidth}\raggedright
Dispatch count
\end{minipage} \\
\midrule\noalign{}
50 & NONE & 0\% \(\pm\) 0.0 & 50.0 \(\pm\) 0.0 \\
50 & BASIC & 32.8\% \(\pm\) 4.1\% & 33.6 \(\pm\) 2.1 \\
50 & ADVANCED & 32.8\% \(\pm\) 4.1\% & 19.0 \(\pm\) 4.1 \\
50 & AGGRESSIVE & 32.8\% \(\pm\) 4.1\% & 19.0 \(\pm\) 4.1 \\
200 & NONE & 0\% \(\pm\) 0.0 & 200.0 \(\pm\) 0.0 \\
200 & BASIC & 34.2\% \(\pm\) 1.8\% & 131.6 \(\pm\) 3.6 \\
200 & ADVANCED & 34.2\% \(\pm\) 1.8\% & 74.0 \(\pm\) 6.3 \\
200 & AGGRESSIVE & 34.2\% \(\pm\) 1.8\% & 74.0 \(\pm\) 6.3 \\
1,000 & NONE & 0\% \(\pm\) 0.0 & 1,000.0 \(\pm\) 0.0 \\
1,000 & BASIC & 35.4\% \(\pm\) 1.6\% & 646.0 \(\pm\) 16.4 \\
1,000 & ADVANCED & 35.4\% \(\pm\) 1.6\% & 344.6 \(\pm\) 31.6 \\
1,000 & AGGRESSIVE & 35.4\% \(\pm\) 1.6\% & 344.6 \(\pm\) 31.6 \\
\bottomrule\noalign{}
\end{tabular}
\end{table}

\textbf{Interpretation.} Three observations:

\begin{enumerate}
\def\labelenumi{\arabic{enumi}.}
\item
  \textbf{Effective instruction count reduction is level-stable}: BASIC,
  ADVANCED, and AGGRESSIVE produce identical reductions (34.2\% \(\pm\)
  1.8\% at length 200). BASIC's fusion and DCE already capture the
  available elimination gains; higher levels add \emph{scheduling}
  improvements (dispatch folding), not elimination---ADVANCED and
  AGGRESSIVE coincide here, their extra passes finding no further
  targets on these sequences, differing only on streams carrying
  combinable patterns. We report this as-is rather than claim monotonic
  gains.
\item
  \textbf{Dispatch count drops further with parallel grouping}: at
  length 200, dispatch falls from 131.6 (BASIC) to 74.0 (ADVANCED) with
  effective-instruction reduction unchanged---adjacent side-effect-free
  instructions with disjoint dependencies and outputs grouped into
  shared dispatches. The 74.0 figure is a scheduling reduction, not a
  further program-size reduction. (An earlier implementation folded more
  aggressively---44.4 at length 200---but could place side-effectful
  instructions into shared groups whose non-interference was never
  established; the sounder grouping here folds less, and Appendix C.5's
  dispatch-folding corollary now matches the implementation exactly.)
\item
  \textbf{Gains are stable across sequence lengths}: 32.8\%, 34.2\%,
  35.4\% at lengths 50, 200, 1,000 respectively, tightest variance at
  length 1,000 (\(\pm\) 1.6\%); length does not materially change the
  opportunity structure.
\end{enumerate}

Effective-instruction reduction is again a \emph{program-size} metric:
it translates to simulated API-call savings only where eliminated or
fused instructions carry model inference; the caching experiment
measures the complementary reuse effect.

\subsection{Experiment 3: Dynamic Hot-Path
Optimization}\label{experiment-3-dynamic-hot-path-optimization}

\textbf{Research question.} Does the dynamic optimizer detect and
specialize repeatedly executed traces?

We record an instruction trace of length 100 for up to 200 iterations
per seed with hot-path detection active (Section 5.1.2). Once a path is
flagged hot, the runtime forms an optimized trace; \textbf{this
experiment measures that trace's length and re-executes it}---the replay
verification below checks that it reproduces the original trace's final
state exactly. Whether a hot-path-optimized stream preserves a
task-level artifact is established separately, on realistic streams, by
Section 5.8 (byte-identical output).

\begin{table}[t]
\caption{Hot-path trace shrinkage per seed (length-100 traces, detection
threshold 50 recordings). The replay column records the executed-state
identity check: the optimized trace, re-executed against an independent
state, reproduces the original trace's final state exactly in every
seed.}
\small\setlength{\tabcolsep}{3pt}
\begin{tabular}{@{}
  >{\raggedright\arraybackslash}p{(\linewidth - 8\tabcolsep) * \real{0.0580}}
  >{\raggedright\arraybackslash}p{(\linewidth - 8\tabcolsep) * \real{0.2174}}
  >{\raggedright\arraybackslash}p{(\linewidth - 8\tabcolsep) * \real{0.2319}}
  >{\raggedright\arraybackslash}p{(\linewidth - 8\tabcolsep) * \real{0.1739}}
  >{\raggedright\arraybackslash}p{(\linewidth - 8\tabcolsep) * \real{0.3188}}@{}}
\toprule\noalign{}
\begin{minipage}[b]{\linewidth}\raggedright
Seed
\end{minipage} & \begin{minipage}[b]{\linewidth}\raggedright
Original length
\end{minipage} & \begin{minipage}[b]{\linewidth}\raggedright
Optimized length
\end{minipage} & \begin{minipage}[b]{\linewidth}\raggedright
Shrink ratio
\end{minipage} & \begin{minipage}[b]{\linewidth}\raggedright
Replay state identical
\end{minipage} \\
\midrule\noalign{}
0 & 100 & 30 & 70.0\% & yes \\
1 & 100 & 41 & 59.0\% & yes \\
2 & 100 & 30 & 70.0\% & yes \\
3 & 100 & 36 & 64.0\% & yes \\
4 & 100 & 34 & 66.0\% & yes \\
\bottomrule\noalign{}
\end{tabular}
\end{table}

Across seeds, hot-path optimization shrinks the trace by \textbf{65.8\%
\(\pm\) 4.6\%}; the seed-to-seed spread (59.0\%--70.0\%) reflects which
sub-sequences each synthetic trace makes hot, the effect present in all
five seeds. As with the static experiments, this is trace-length
reduction on synthetic traces---workload-controlled evidence that the
detection-formation-reoptimization loop functions, not a wall-clock
speedup claim.

\textbf{Replay verification.} Each seed's optimized trace \emph{is}
re-executed---every instruction reads its declared dependencies from a
threaded state dict and returns deterministic bindings for its declared
outputs (side-effectful instructions leave per-id markers)---and its
final state is compared with the original's, executed the same way. The
states are identical in all five seeds (5/5; recorded per seed as
\texttt{replay\_state\_identical} in the results file), so the 65.8\%
shrinkage carries an executed-equality check at trace level---on the
deterministic mock, not against a live model; this is the same style of
check as Section 5.8's task-level byte-identity, applied to the
synthetic traces. Live-model behavior remains untested---no mock
experiment can address it.

\subsection{Ablation: Per-pass
Contribution}\label{ablation-per-pass-contribution}

To attribute the static gains we run the two ablations of Section 5.1.2.
The primary one is \textbf{paired}: it compiles the \emph{same}
sequences as Section 5.3 (identical seed base 2,000 and generator) with
fusion and DCE toggled individually and jointly, so per-configuration
differences are same-input paired differences---an exact decomposition
of the headline, not a replication:

\begin{table}[t]
\caption{Paired per-pass decomposition on the Section 5.3 sequences
(length 200, 5 seeds, identical seed base 2,000); the fusion + DCE row
reproduces the Section 5.3 headline exactly because the inputs are
identical.}
\small\setlength{\tabcolsep}{3pt}
\begin{tabular}{@{}ll@{}}
\toprule\noalign{}
Configuration & Instruction-count reduction \\
\midrule\noalign{}
NONE & 0.0\% \(\pm\) 0.0\% \\
fusion only & 21.6\% \(\pm\) 1.2\% \\
DCE only & 24.5\% \(\pm\) 2.6\% \\
fusion + DCE & 34.2\% \(\pm\) 1.8\% \\
\bottomrule\noalign{}
\end{tabular}
\end{table}

The fusion+DCE row is 34.2\% \(\pm\) 1.8\%---the headline itself,
confirming that these two passes account for the entire static gain at
length 200. The \textbf{paired interaction}, computed per seed as
(fusion+DCE − NONE) − (fusion − NONE) − (DCE − NONE)---34.2 − 21.6 −
24.5 = −11.9 points---is −11.9\% \(\pm\) 2.0\% (per-seed
−14.5/−9.5/−13.0/−12.0/−10.5 points): the passes overlap---dead slots
break fusion chains and fused pairs absorb neighbors DCE would otherwise
have removed individually---so the sum of the single-pass gains (46.1
points) overstates the joint effect by 11.9 points.

The secondary ablation is the earlier \textbf{independent-draw
replication} (seed base 4,000, same distribution): fusion alone 21.9\%
\(\pm\) 3.5\%, DCE alone 24.1\% \(\pm\) 3.3\%, fusion+DCE and all passes
35.0\% \(\pm\) 3.7\%---the paired structure reproduces on fresh
sequences (single-pass shares agree within one standard deviation), so
the decomposition is not an artifact of one draw. ``All passes'' equals
fusion+DCE here, consistent with Section 5.3's level stability: no
additional pass contributes measurable reduction, and we report that
absence rather than dress the higher levels in unearned credit. Because
dead instructions are planted at rate 0.25 by construction, DCE's
per-pass share is near-tautological for this family---it measures that
the pass finds the planted dead code, precisely its function.

\subsection{Multi-Agent Coordination: Functional
Validation}\label{multi-agent-coordination-functional-validation}

The multi-agent layer---\texttt{AgentCoordinator},
\texttt{SecurityMonitor}, and the collaborative instruction set---is
validated functionally by 43 tests, all passing---the
coordination-and-security subset of the repository's full suite (861
passing, 14 skipped at run time), which also covers the ISA core,
compiler, and optimizer units of Sections 5.2--5.5. This is
\textbf{qualitative functional validation}: the tests verify the
coordination behaviors exist and are correct, not that they are fast; no
performance number should be inferred here.

The coordination flow implemented by \texttt{AgentCoordinator}
(authenticate, authorize, select by capability, plan stages with
fallback, execute, record history) is listed as pseudocode in Appendix
D.10 of the supplementary material. What the 43 tests verify, grouped by
subsystem:

\begin{enumerate}
\def\labelenumi{\arabic{enumi}.}
\tightlist
\item
  \textbf{Agent registration and role policies}: registration creates
  policies; permissions and security requirements differ by role; trust
  levels are bounded and reward declared specializations.
\item
  \textbf{Task coordination}: coordination produces a solution or
  records failure; missing authentication/authorization and
  no-suitable-agents raise; selection respects the capability threshold,
  the coordinator being the highest-scoring agent.
\item
  \textbf{Execution planning}: stage structure with coordination rules;
  every stage has differentiated fallback keys, and a fallback executes
  on stage failure; validators are reasoning-heavy non-coordinator
  agents.
\item
  \textbf{SecurityMonitor}: valid tasks pass, missing
  authentication/authorization fails, violations are timestamped,
  declared data sensitivity is enforced (sensitive categories infer a
  high sensitivity level, benign ones pass at medium), resource access
  is open by default and restrictable by configuration.
\item
  \textbf{Collaborative instructions}: barriers wait for the declared
  participant count and time out on a missing peer; resource allocation
  rejects over-allocation and serializes concurrent requests; the
  workflow state machine accepts valid actions and rejects invalid ones
  and double starts; a research-team scenario runs end to end.
\end{enumerate}

These tests exercise the same production modules used elsewhere and are
fully deterministic (no LLM calls). They establish correctness of the
coordination machinery under the tested conditions; what they do not
establish is scalability under contention, crash recovery, or
adversarial inputs beyond the policy checks above---all future work
(Section 6).

\subsection{Real-LLM End-to-End Validation of Result
Caching}\label{real-llm-end-to-end-validation-of-result-caching}

\textbf{Research question.} Does the caching mechanism behave
identically when the ``model'' is a real LLM behind a live gateway, and
how much \emph{measured} wall-clock time and token usage does it save?

\textbf{Setup.} We replace the mock backend with a real model behind a
local Anthropic-protocol gateway (TLS with an internal CA; alias
\texttt{claude-glm-5.3-flash}, a hosted reasoning-capable GLM model;
\texttt{max\_tokens=64}; sequential httpx, 90 s timeout, three bounded
retries---none occurred). Every other component is the production code
path: the same \texttt{InstructionCache} protocol, instruction objects,
and workload generator, with Section 5.2's instruction pool (50
templates) and parameter space (20 values).

\textbf{Scale.} With each round trip costing \textasciitilde2.2 s, scale
drops from 2,000 tasks \(\times\) 5 seeds to \textbf{200 tasks
\(\times\) 1 seed} per Zipf level (s in \{0.5, 1.0, 1.5\}): 1,040
measured round trips (plus one preflight) in 39.4 minutes. The seed pins
the task sequence exactly, so the workload is the one the expectation
below is computed from.

\textbf{An exactness property.} A cache hit depends only on the
(instruction, params) pairs seen so far---not on what the model returns.
The production \texttt{InstructionCache} is LRU with a TTL, so in
general a hit also depends on wall-clock time and capacity; under this
benchmark none of that can bind: \texttt{max\_size} = 100,000 (at most
200 distinct keys per arm), \texttt{ttl} = 3,600 s (every cache-enabled
arm completes in under seven minutes; the no-cache arms run to 7.9
minutes and never consult a cache), a fresh cache per arm, failed calls
never cached. Hit/miss decisions are therefore a pure function of the
(instruction, params) sequence, \emph{computable without any backend};
if the mechanism is correct the real run must reproduce that sequence
exactly---and at the recorded granularity (per-arm hit rates and call
counts) it does:

\begin{table}[t]
\caption{Real-gateway cache behavior against the computed expectation
(200 sequential requests, single seed, three Zipf levels, single model;
paired cells read no cache -\textgreater{} cached).}
\small\setlength{\tabcolsep}{3pt}
\begin{tabular}{@{}
  >{\raggedright\arraybackslash}p{(\linewidth - 14\tabcolsep) * \real{0.0359}}
  >{\raggedright\arraybackslash}p{(\linewidth - 14\tabcolsep) * \real{0.2395}}
  >{\raggedright\arraybackslash}p{(\linewidth - 14\tabcolsep) * \real{0.1018}}
  >{\raggedright\arraybackslash}p{(\linewidth - 14\tabcolsep) * \real{0.1617}}
  >{\raggedright\arraybackslash}p{(\linewidth - 14\tabcolsep) * \real{0.0838}}
  >{\raggedright\arraybackslash}p{(\linewidth - 14\tabcolsep) * \real{0.2156}}
  >{\raggedright\arraybackslash}p{(\linewidth - 14\tabcolsep) * \real{0.0898}}
  >{\raggedright\arraybackslash}p{(\linewidth - 14\tabcolsep) * \real{0.0719}}@{}}
\toprule\noalign{}
\begin{minipage}[b]{\linewidth}\raggedright
Zipf s
\end{minipage} & \begin{minipage}[b]{\linewidth}\raggedright
Expected hit rate (computed, no backend)
\end{minipage} & \begin{minipage}[b]{\linewidth}\raggedright
Observed hit rate
\end{minipage} & \begin{minipage}[b]{\linewidth}\raggedright
Calls
\end{minipage} & \begin{minipage}[b]{\linewidth}\raggedright
Call reduction
\end{minipage} & \begin{minipage}[b]{\linewidth}\raggedright
Wall clock (ms)
\end{minipage} & \begin{minipage}[b]{\linewidth}\raggedright
Wall saved (ms)
\end{minipage} & \begin{minipage}[b]{\linewidth}\raggedright
Tokens saved
\end{minipage} \\
\midrule\noalign{}
0.5 & 8.5\% & 8.5\% & 200 -\textgreater{} 183 & 8.5\% & 472,276
-\textgreater{} 405,754 & 66,522 & 1,445 \\
1.0 & 26.0\% & 26.0\% & 200 -\textgreater{} 148 & 26.0\% & 455,946
-\textgreater{} 334,597 & 121,349 & 4,420 \\
1.5 & 45.5\% & 45.5\% & 200 -\textgreater{} 109 & 45.5\% & 450,330
-\textgreater{} 244,316 & 206,014 & 7,735 \\
\bottomrule\noalign{}
\end{tabular}
\end{table}

Per-call latency was stable across arms (p50 2.13--2.21 s, p95
3.06--3.54 s). Token savings are gateway-reported \texttt{usage}
differences: every avoided call saves its input tokens plus the tokens
the reasoning model generates before responding (output tokens were the
binding constraint of \texttt{max\_tokens=64} in every call).

\textbf{Interpretation.} First, the mechanism transfers unchanged:
observed hit rates match the computed expectation exactly at every skew
level, so Section 5.2's synthetic results measure the cache, not a mock
artifact. Second, avoided calls convert into measured time and
tokens---14.1\%, 26.6\%, and 45.8\% of the no-cache arm's wall clock at
s = 0.5, 1.0, 1.5---tracking but not identical to the call reduction,
since arms ran sequentially and per-call latency drifted slightly (the
exact quantity is the call reduction; wall and token figures are
measurements, not derivations). Third, the \emph{absolute} hit rates
here (8.5--45.5\%) are lower than Section 5.2's (59.5--79.6\%) purely
because repetition accumulates with task count: at N = 200 these are the
exact computed values for this workload family, and at N = 2,000 the
same computation yields the Section 5.2 numbers---a property of the
workload, not the backend.

\textbf{Scope of the claim.} This experiment validates the caching
mechanism end-to-end---one gateway, one model, one seed, N = 200,
sequential requests, no streaming---recording per-arm aggregate hit
rates and call counts; per-position identity follows from the
deterministic cache decisions above but is not separately logged, and
nothing here is a throughput, availability, or dollar-cost claim.
Sections 5.3--5.5 are structural, never dispatching to the backend;
Section 5.7.1 extends the real-gateway validation to them via a
cache-mediated protocol, their validity otherwise resting on the
code-path and ablation arguments there.

\subsubsection{Real-Gateway Validation of Compilation and Hot-Path
Optimization}\label{real-gateway-validation-of-compilation-and-hot-path-optimization}

\textbf{Research question.} Do static compilation and closed-loop
hot-path adoption preserve the observable semantics of an instruction
stream when the stream's LLM step is served by a real model---and what
do the mechanisms measurably save on a live gateway?

\textbf{Setup, arms, scale.} A dedicated harness
(\texttt{examples/paper\_real\_llm\_compiler\_test.py}; tests in
\texttt{tests/system/test\_real\_llm\_compiler\_test.py}) reuses Section
5.7's gateway configuration (same base URL, internal-CA TLS, alias
\texttt{claude-glm-5.3-flash}, sequential httpx, 90 s timeout, three
bounded retries) but raises \texttt{max\_tokens} to 512 (Section 5.7
used 64) so completions dominate the priced comparison; it runs
identically online or offline against a deterministic same-protocol mock
(call counting, template-accounted tokens, constant 80 ms simulated
latency). The task is Section 5.8's service-log triage shape: a
seven-instruction stream with one cacheable LLM-type step
(\texttt{classify\_signature}, one model call per distinct
incident-signature batch) and the same planted dead guard
(\texttt{check\_format}), over 50 realizations \(\times\) 1 seed
\(\times\) Zipf s = 1.0 from a fixed universe of 40 task instances---22
distinct signature batches, so the entire run is 23 gateway round trips
(one preflight plus 22 live calls, all by arm (a)) in 4.4 minutes, a
bounded validation run, not the 200 \(\times\) 3-skew sweep of Section
5.7. Arm (a) is the unoptimized authored stream, run first, cold,
performing every live round trip; arm (b) drives the \emph{same} stream
through the production machinery---\texttt{optimize\_hot\_path} adoption
plus static \texttt{InstructionCompiler} at NONE vs AGGRESSIVE with the
dataflow-preservation assertion---then executes the same draws against
arm (a)'s warmed cache. The classify step's cache key is the sha256 of
exactly the input dictionary its prompt is built from, so arm (b)
issuing \emph{zero} new gateway calls proves it computed identical
keys---hence identical prompts and consumed values---and its artifacts
must come out byte-for-byte equal; any perturbation of the LLM step's
inputs would surface as a cache miss and/or artifact divergence.

\textbf{Results.} All checks pass. Arm (a) made exactly the 22
first-occurrence calls the backend-free computation predicts and served
the other 28 realizations from the result cache (hit rate 56.0\%),
consuming 4,027 prompt and 11,264 gateway-reported completion tokens in
260.2 s of measured call latency. Arm (b)---compiled,
hot-path-adopted---made zero new gateway calls and reproduced every
artifact byte-for-byte. The hot path was adopted at realization 7
(threshold 5; adoption lands two realizations past it because the
counter must strictly exceed it and recording follows execution, as in
Section 5.8): realizations 1--6 executed the full seven-instruction
stream, realizations 7--50 the shrunken six-instruction form with
\texttt{check\_format} absorbed into its successor by fusion---306
instruction executions against 350 in arm (a), a 12.6\%
executed-instruction reduction attributable to the adopted trace
(per-stream shrink 14.3\%, 7 \(\rightarrow\) 6; realized saving diluted
by the six full-form realizations before adoption; the static compiler
alone reproduces the 7 \(\rightarrow\) 6 shrink at AGGRESSIVE, so the
two mechanisms' contributions are not separated), reported separately
from the cache's call savings. Static compilation of the same stream:
NONE 7 \(\rightarrow\) AGGRESSIVE 6 effective instructions, dataflow
preserved.

\textbf{A priced comparison (summary; protocol and table in Appendix G
of the supplementary material).} Costs are MODELED---computed from a
pricing overlay on measured tokens, not observed: the overlay prices the
run's token counts under no optimization, AgentISA result caching (first
occurrences only---the measured behavior of both arms), and a modeled
provider prompt cache (every realization queries; repeat prompts bill
input tokens as prompt-cache reads at 0.1x, mirroring Anthropic's
published read discount; completions always bill at full price). Two
simplifications---omitting Anthropic's 1.25x cache-write premium and
1024-token minimum cacheable prefix---both favor the provider-cache
baseline, so the comparison is conservative for us; list prices are
parameterized stand-ins (\$3/M input, \$15/M output) because the
gateway's actual pricing is unknown; no dollar-savings claim is made. On
this run's token mix the result cache cuts modeled cost by 56.0\%---it
removes whole calls, completions included---while the modeled provider
prompt cache cuts only 3.3\%, the reasoning model's completions
dominating; on the deterministic mock at paper scale (200 tasks
\(\times\) 1 seed \(\times\) 3 Zipf levels) the result cache cuts
modeled cost by 80.1\%, 81.6\%, and 85.1\% at s = 0.5, 1.0, 1.5, against
28.9--31.4\% for the provider prompt cache---there short completions do
not dominate, so the input-side discount matters more; the ranking of
the two caching strategies is a function of the input/output token mix,
not a constant. GPTCache, an external cache layer, is explicitly out of
scope; no head-to-head against it is claimed.

\textbf{What this does and does not show.} The byte-identity property is
cache-mediated by design: the cached values entered the artifacts once,
through arm (a)---values that on this alias are fallback-derived, as
explained below---and arm (b) is proven to compute the same keys and
consume the same values; this validates the stream transformation, not
model determinism, which is not guaranteed at temperature 0 and is not
claimed. On this alias the reasoning model's replies bind at
\texttt{max\_tokens} in every call (as at \texttt{max\_tokens=64} in
Section 5.7, and as a 1,500-token diagnostic still did), so the classify
step's severity values come from its deterministic fallback; what is
validated against the live model is the mechanism chain (key
computation, cache decisions, call and token accounting, artifact
identity), not classification quality.

\textbf{What the real-gateway experiments do and do not establish.}
Three layers, kept separate: (1) \emph{gateway integration
correctness}---requests reach the gateway, the cache governs calls, and
call/token statistics are recorded faithfully: \textbf{established} by
Sections 5.7 and 5.7.1. (2) \emph{Workflow semantic preservation}---the
optimized stream computes the same cache keys and consumes the same
cached values as the unoptimized one, yielding byte-identical artifacts:
\textbf{established as a cache-mediated identity} (see above),
validating the stream transformation, not model determinism. (3)
\emph{Real-model task quality}: \textbf{not established here}---the
reasoning alias returned results at the token ceiling and the classify
step fell back to a deterministic rule, so no statement about live-model
output quality is made anywhere in this paper; the abstract, highlights,
and conclusion state only layers (1) and (2). Both real-gateway runs
share the coverage limits above (one gateway, one model, one seed,
sequential, no streaming); multi-model, multi-seed, and real-deployment
behavior remain untested.

\subsection{Controlled Four-Task System-Level
Characterization}\label{controlled-four-task-system-level-characterization}

The experiments above measure individual mechanisms on synthetic
sequences; this section is a \emph{controlled four-task system-level
characterization}---a functional demonstration under a fixed closed-loop
protocol (same task, 12 realizations, one cacheable step each), not a
generalization experiment; the identical 91.7\% hit-rate figure below is
a property of that protocol, not four independent cache observations.
The seed-controlled characterization
(\texttt{examples/system\_task\_workload.py},
\texttt{tests/system/test\_system\_task\_workload.py}; the original
single-task demonstration lives on as
\texttt{examples/system\_task\_demo.py}) is built from \textbf{four
tasks with genuinely different dependency shapes}---linear service-log
triage, fan-out csv-etl-anomaly, wide-fan-in report-merge, two-branch
regression-diff---each realized end-to-end against sandbox-resident
files (12 realizations each, hot-path threshold 5, one planted dead
guard and one cacheable LLM-type step per task, the latter backed by a
deterministic rule stand-in counting simulated model calls), making the
mechanisms' composition---and Section 3.4.4's closed-loop optimization
claims---checkable. Per-task instruction streams, the single-task
walk-through, and both result tables are in Appendix H of the
supplementary material; the headline: every task's guard survives
compilation (never dropped outright, the reduction being fusion), the
first realization pays all model calls and every subsequent one hits the
cache (91.7\%, 11 of 12 realizations by construction---one cacheable
step, identical seeded draws), hot-path adoption lands at repeat 7
executing the shrunken six-instruction stream, and all reports come out
byte-identical (12/12 plus fresh-stack SHA-256 equality); the compile
shrink is the only shape-dependent quantity (12.5--25.0\%), which is
what one wants from an infrastructure claim---the composition behavior
is not an artifact of one stream's structure.

\textbf{The triage task's instruction stream} (seven instructions over
sandbox-resident files): \texttt{ingest\_logs} reads six 400-line
sandbox log files (side-effectful); \texttt{check\_syntax} is a
defensive format guard that never rejects here (a \emph{planted dead
guard}, standing in for real pipelines' belt-and-braces checks);
\texttt{normalize\_entries} parses and normalizes lines;
\texttt{extract\_incidents} lifts them into incident records;
\texttt{classify\_signature} triage-classifies each incident (LLM-type,
cacheable, backed by a deterministic call-counting rule stand-in);
\texttt{aggregate\_triage} summarizes by severity;
\texttt{write\_report} writes the report (side-effectful).

\textbf{Calibrating the synthetic headline.} Section 5.4 reports 65.8\%
trace shrinkage on synthetic traces with planted dead instructions; the
four realistic tasks show \textbf{12.5--25.0\%} (one planted guard in a
seven- or eight-instruction stream, plus one additional live-pure fusion
pair in csv-etl-anomaly). We state both figures rather than privilege
either: 12.5--25.0\% reflects the planted dead-check density (one guard
per seven- to eight-instruction stream)---a stated density assumption,
not observed ordinary-pipeline waste---while the synthetic figure is an
upper bound under the skew/planting parameters we control, not a
typical-case estimate. Constructing a realistic (non-synthetic)
dependency shape also served the audit purpose, exposing two
grouping/fusion soundness gaps---side-effectful instructions opening
dispatch groups, fused instructions losing their external
dependencies---now fixed with regression tests; Section 5.3's
dispatch-folding numbers reflect the corrected (more conservative)
grouping.

\textbf{Structural invariance.} Every task passes the same structural
assertions: the guard is never dropped outright (the eliminated set is
empty in all four), the reduction being fusion---in csv-etl-anomaly a
second adjacent pure pair also fuses, hence its 25.0\% shrink with a
single planted guard. The first realization pays all model calls and
every subsequent one hits the cache---the 91.7\% hit rate is 11 of 12
realizations by construction (one cacheable step, identical seeded
draws), so the four tasks share it structurally. Reuse and closed-loop
adoption are thus invariant across the four dependency shapes under the
fixed protocol, while the guard/fusion shrink varies with stream
structure (12.5--25.0\%)---the compile shrink is the only
shape-dependent quantity; the characterization's value over the
single-task demonstration is evidence the composition behavior is not an
artifact of one stream's structure.

\subsection{Limitations and Threats to
Validity}\label{limitations-and-threats-to-validity}

We enumerate the limitations of the evaluation explicitly:

\begin{enumerate}
\def\labelenumi{\arabic{enumi}.}
\item
  \textbf{Mock backend, not a live LLM.} The mock charges a constant 80
  ms per simulated call and returns deterministic responses; beyond
  counting, it covers Section 5.4's replay verification (final-state
  identity, 5/5 seeds) and Section 5.2's naive-cache baseline (same
  instrumentation). The real-gateway experiments cover the cache
  reproducing the computed hit/miss statistics and converting avoided
  calls into measured time and tokens (Section 5.7), and the compiled,
  hot-path-adopted stream issuing zero new gateway calls with
  byte-identical artifacts (Section 5.7.1)---a cache-mediated identity
  validating the stream transformation, not model determinism. Still not
  covered: multi-model and multi-seed behavior, real deployment, caching
  under model nondeterminism \emph{across prompt variations}, streaming
  responses, provider rate limits, and live-model task quality (the
  four-task classification steps are deterministic rule-based stand-ins,
  and the 5.7.1 gateway run fell back to the deterministic classifier).
\item
  \textbf{Synthetic workloads.} Zipf-distributed synthetic tasks are a
  reasonable model of skewed agent workloads, but a model; Section 5.8's
  task-level characterization only partially addresses this. The
  cache-reduction figures (59.5--79.6\%) are exact for these workload
  families and should be read as \emph{skew-sensitivity curves}, not
  universal savings; uniformly novel tasks would see near-zero benefit.
  Section 5.8's four-shape characterization calibrates the 65.8\%
  synthetic headline downward with its 12.5--25.0\% per-task compile
  shrink and shows uniform reuse and closed-loop adoption under the
  fixed protocol (12 seeded realizations, threshold 5)---but the four
  tasks remain seed-controlled, hand-constructed streams, not a
  production workload study.
\item
  \textbf{No production end-to-end evaluation.} We have not deployed
  AgentISA in a production setting; consequently this paper contains no
  latency-under-load, throughput, availability, or dollar-cost
  measurements, and we make no such claims.
\item
  \textbf{No real cost data (token usage measured; pricing MODELED on
  measured tokens).} All latency figures in Sections 5.2--5.5 are
  simulated (avoided calls \(\times\) 80 ms); Section 5.7 measures real
  wall-clock time and gateway-reported tokens at reduced scale, and
  Section 5.7.1 prices its own measured tokens (the 22-call reference
  run) under three modes---no optimization, AgentISA result caching, and
  a modeled provider prompt cache with a parameterized cache-read
  discount (default 0.1, after Anthropic's published prompt-caching
  economics). The costs are MODELED-on-measured: tokens are
  measurements, prices are configurable list-price stand-ins (the
  gateway's actual pricing is unknown), and no dollar-savings claim is
  made. Section 5.2's baseline is an \emph{exact-match} response cache
  (parity reported); head-to-head against \emph{approximate}
  (semantic/embedding) caching (GPTCache's approximate mode---explicitly
  out of scope) and a priced deployment study under real pricing remain
  future work.
\item
  \textbf{Program-size metrics, not wall-clock time.}
  Effective-instruction and dispatch-count reduction are structural
  metrics: they correlate with---but do not equal---wall-clock
  improvements, which depend on per-instruction cost, available
  parallelism, and dispatch overhead.
\item
  \textbf{Descriptive statistics only.} Synthetic configurations use
  five seeds (Section 5.1.3); the real-gateway and system-level runs are
  single-seed by design. No hypothesis tests, confidence intervals, or
  significance claims are made; standard deviations indicate stability
  across seeds, nothing more.
\item
  \textbf{Functional, not performance, validation of coordination.} The
  43 multi-agent tests verify behavioral correctness under deterministic
  conditions, not throughput or scalability.
\end{enumerate}

Despite these limitations, the evaluation achieves its purpose: it
demonstrates, under fully disclosed and reproducible conditions, that
each optimization mechanism is implemented, active, and produces the
claimed structural effects on workloads designed to expose them.

\section{Discussion}\label{discussion}

\subsection{System Limitations and
Challenges}\label{system-limitations-and-challenges}

\begin{enumerate}
\def\labelenumi{\arabic{enumi}.}
\item
  \textbf{Instruction-set coverage}: complex reasoning patterns are not
  yet naturally expressible (domain-specific extensions are likely
  needed), and maintaining consistency of pre/post-conditions and typing
  rules across extensions remains an open engineering challenge.
\item
  \textbf{Evaluation scope}: as detailed in Section 5.9, the
  quantitative evidence is primarily microbenchmark-based (Sections
  5.2--5.5), supplemented by real-gateway spot checks (Sections 5.7 and
  5.7.1) and a seed-controlled four-task characterization (Section 5.8);
  production end-to-end evaluation---including real provider pricing,
  rate limits, and failure modes---is required before any deployment
  claim.
\item
  \textbf{Performance bottlenecks}: cache memory consumption under large
  instruction universes; synchronization overhead in barrier-coordinated
  multi-agent execution; recompilation cost when hot-path thresholds are
  misconfigured; the shipped hot-path detector keys on exact sequence
  identity, so trace diversity never accumulates counts (the EMA-based
  windowed detector specified in Section 4.2.2 is the designed remedy);
  hot-path state is currently unbounded---production use needs an
  eviction policy and a cap on tracked paths.
\item
  \textbf{Integration challenges}: adapter maintenance as provider
  function-calling surfaces evolve; version management of learned
  instructions across LLM updates.
\end{enumerate}

\subsection{Security and Privacy
Considerations}\label{security-and-privacy-considerations}

AgentISA's implemented security layer is the \texttt{SecurityMonitor} of
Section 5.6---task-level authentication, role-based authorization,
data-sensitivity classification, configurable resource-access
policies---consulted by the coordinator before every assignment, with
violations timestamped. One shipped default deserves an explicit
call-out: \textbf{resource access is open by default} and restrictable
by configuration---a fail-open stance shifting the security burden onto
the policy writer; a hardened deployment should invert it to fail-closed
(deny on unregistered resources, unknown sensitivity, or incomplete
policy), which the current checks do not do. The 43 functional tests
validate the implemented checks under configured policies, not the
default posture; nothing in this paper constitutes a security
certification. Beyond what is implemented, hardening
considerations---not claims about the current system: (i) \textbf{data
privacy}---end-to-end data-flow tracking and privacy-preserving
computation remain future work (Section 6.3.6); (ii) \textbf{system
security}---secure multi-agent communication channels and audit-log
integrity are not yet addressed; (iii) \textbf{compliance}---regulated
deployments require externally audited controls; the system provides
enforcement hooks, not certification.

\subsection{Future Research
Directions}\label{future-research-directions}

Readers pressed for time need only the first tier. \textbf{Directly
motivated by this paper's exposed gaps:} 6.3.1 (unifying the formal
semantics and discharging the conditional theorems---responding to
Sections 3.3.3 and C.7's specified-vs-shipped split), 6.3.4 (calibrating
capability scores against end-to-end performance---responding to Section
3.5's stub disclosure), and the real-workload evaluation named in
6.3.3's first item (responding to Section 5.9). \textbf{Longer-term,
with no current implementation link:} 6.3.2 (self-modifying and
context-aware instruction sets---formulas are sketches, as marked
inline), 6.3.5 (emergence metrics, cross-agent transfer), and the
privacy composition of 6.3.6.

\subsubsection{Theoretical Foundations}\label{theoretical-foundations}

\begin{enumerate}
\def\labelenumi{\arabic{enumi}.}
\item
  \textbf{Formal semantics of agent instructions}: consolidating Section
  3.1 and Appendix B into one standalone semantics---instruction
  language, state space, reduction rules, type system---is the
  foundational next step; nothing in it is new machinery, and the
  composition is what a journal-length treatment would spell out.
\item
  \textbf{Compositional reasoning}: a parallel-composition rule whose
  disjointness side condition mirrors the \(\parallel\)-typing rule of
  Definition 3.6---without it the conjunction of preconditions does not
  guarantee either component's precondition during interleaving---would
  enable modular verification, parallel composition, and temporal
  property preservation (rule stated in Appendix E.1 of the
  supplementary material).
\item
  \textbf{Type theory for agent capabilities}: a capability type
  discipline with Basic/Composite/Emergent kinds and a composition rule
  \(a \circ b : \text{cap}_1 \otimes \text{cap}_2\) (Appendix E.1).
\end{enumerate}

\subsubsection{Advanced Instruction Set
Evolution}\label{advanced-instruction-set-evolution}

\begin{enumerate}
\def\labelenumi{\arabic{enumi}.}
\item
  \textbf{Self-modifying instructions}---a formal model for instruction
  evolution (f is left abstract; the sketch names the inputs, not a
  concrete update rule):

  \[
  I_{t+1} = f(I_t, \text{Performance}_t, \text{Context}_t)
  \]

  under semantic-preservation, performance-monotonicity, and safety
  constraints.
\item
  \textbf{Context-aware optimization}---for each observed context,
  predict the instruction set's performance against a threshold, and if
  it falls short, select a best rewrite variant for that context and
  accumulate it into the set. The programmatic sketch (OPTIMIZE-CONTEXT)
  is listed in Appendix D.9; its primitives are future work, not
  implemented machinery.
\item
  \textbf{Cross-model instruction mapping}: a universal mapping
  framework \(\phi_{A \rightarrow B}: I_A \rightarrow I_B\) with
  semantic preservation, performance bounds, and capability awareness
  stated as goals, not delivered guarantees.
\end{enumerate}

\subsubsection{Theoretical Performance
Bounds}\label{theoretical-performance-bounds}

\begin{enumerate}
\def\labelenumi{\arabic{enumi}.}
\item
  \textbf{Asymptotic analysis}: deriving cost bounds per composition
  operator---when does iterative composition \(i_1^*\) admit a
  closed-form cost model, and under what reuse structure does caching
  yield amortized \(O(1)\) effective cost per task? The microbenchmarks
  (Section 5) measure the constants empirically; the asymptotic theory
  remains open.
\item
  \textbf{Resource consumption models}: a calibrated model
  \(R(i) = \alpha C_{cpu}(i) + \beta C_{mem}(i) + \gamma C_{io}(i)\)
  with platform-specific coefficients fitted against real deployments
  rather than assumed (Appendix E.1).
\end{enumerate}

\subsubsection{Capability Assessment:
Extensions}\label{capability-assessment-extensions}

The instruction-level capability assessment subsystem
(\texttt{CapabilityProfile}, Section 3.5) ships in part---aggregation
and the coordinator's threshold selection are implemented and used,
while instruction-level metric extraction remains a specified stub. Two
extensions remain future work: (i) \textbf{benchmark-calibrated
metrics}, correlating static, instruction-derived domain scores with
end-to-end benchmark performance, in the manner of function-calling
evaluations {[}17{]} and agent benchmarks {[}29{]}, so that weights
\(w_k\) are empirically grounded rather than configured by hand; and
(ii) \textbf{cross-model comparison matrices}, extending the per-agent
profile to a pairwise framework \(\Delta(M_1, M_2)\) over capability
dimensions, supporting systematic model selection for a given
instruction workload.

\subsubsection{Emergent Behavior
Analysis}\label{emergent-behavior-analysis}

\begin{enumerate}
\def\labelenumi{\arabic{enumi}.}
\item
  \textbf{Collective intelligence modeling}: the open question is a
  formal metric for when multi-agent composition exceeds the sum of its
  members' capabilities---what emergence would have to mean measurably,
  rather than an appeal to an undefined \(\varphi\) over unioned
  capability sets.
\item
  \textbf{Learning transfer mechanisms}: knowledge transfer across
  agents composes the mining-extraction-validation pipeline of Section
  3.3 rather than introducing new machinery---patterns mined from a
  source agent's traces would be adapted to a target agent's instruction
  set and re-validated before promotion; the validation obligation
  transfers unchanged.
\end{enumerate}

\subsubsection{Security and Privacy
Extensions}\label{security-and-privacy-extensions}

\begin{enumerate}
\def\labelenumi{\arabic{enumi}.}
\item
  \textbf{Information flow control}: the classical noninterference
  property
  \(\text{Secure}(P) \iff \forall s_1, s_2: L(s_1) = L(s_2) \implies L(\llbracket P \rrbracket s_1) = L(\llbracket P \rrbracket s_2)\),
  where \(L\) assigns security levels to states and outputs, stated for
  terminating runs (divergence and channel leakage out of scope;
  Appendix E.1).
\item
  \textbf{Privacy-preserving computation}: a direction, not a delivered
  mechanism---if agent computations over sensitive data were routed
  through per-step differentially private mechanisms, the overall loss
  would follow a composition theorem over the per-mechanism
  \((\epsilon_i, \delta_i)\) parameters. We deliberately state no bound:
  which composition theorem applies, and which concrete mechanisms the
  runtime would expose, are open design choices; any formula written now
  would be decoration rather than a guarantee. The direction is listed
  because the instruction contract is the natural place to carry such
  annotations once the mechanisms exist.
\end{enumerate}

None of the frameworks above is implemented today; Sections 6.3.3 and
6.3.4 are stated precisely enough to be falsifiable, while the remaining
sketches are programmatic, guiding subsequent work from
controlled-setting-validated infrastructure toward a deployable,
formally verified agent platform.

\section{Conclusion}\label{conclusion}

This paper presented AgentISA, an instruction set architecture that
treats LLM agent behavior as a first-class compilable artifact. Our
contributions are:

\begin{enumerate}
\def\labelenumi{\arabic{enumi}.}
\tightlist
\item
  A formal instruction set with compositional semantics, a type system,
  and pre/post-conditions, with sanity results as design rationale
  (Turing completeness as negative certification, Theorem 3.1; pure-core
  type safety, Theorem 3.2; extraction-time context preservation,
  Theorem 3.3) and semantic-preservation arguments for the optimization
  and scheduling transformations (Theorems 3.4--3.6).
\item
  Agentic Calling, which lifts provider-specific function calling into
  reusable, provider-independent, cacheable agent instructions.
\item
  Trace-based instruction learning---frequency mining with deterministic
  action-signature grouping, developer-controlled promotion, and
  composite replay.
\item
  A compiler and runtime applying classical optimizations---fusion,
  dead-code elimination, scheduling, and dynamic hot-path
  specialization---to agent instruction streams.
\item
  A multi-agent coordination layer with capability-based selection and a
  security monitor enforcing authentication, authorization, and
  data-sensitivity policies.
\end{enumerate}

The evaluation shows the optimizations are real and measurable under
controlled conditions: 59.5--79.6\% simulated-call reduction from
caching, 34.2\% instruction-count reduction from static compilation,
65.8\% trace shrinkage from hot-path optimization (Sections 5.2--5.5),
real-gateway validations of the cache and the compiled pipeline
(Sections 5.7 and 5.7.1, gateway-integration and cache-mediated-identity
layers only), a four-task system-level characterization (Section 5.8),
and 43 functional tests for the multi-agent layer. These are
microbenchmark and controlled-setting results, not production end-to-end
measurements; the threats to validity of Section 5.9 stand.

\textbf{Implementation status of the contributions.} For auditability,
the five contributions above split by status: \emph{implemented and
exercised by the benchmarks}---the instruction set and contracts,
Agentic Calling, result caching, fusion/DCE/grouping, hot-path detection
and adoption, the coordination layer with injected-score capability
selection, and the security monitor (43 tests across both);
\emph{implemented with known limits}---the bounded checker (pure
instructions, bounded state spaces only), runtime context validation
(deletion blind spot disclosed), and pattern mining with composite
replay (developer-controlled promotion, no generalization beyond
recorded parameters); \emph{specified, not implemented}---tree-pattern
instruction selection, cycle-accurate pipeline scheduling,
symbolic-equivalence checking, and capability metric extraction.
Theorems 3.4--3.6 are conditional arguments for the implemented passes;
the specified components carry no verification claim.

AgentISA's longer-term significance is architectural: by giving agent
behavior an instruction-level contract, it makes agent systems amenable
to the same discipline that transformed hardware and compiler
engineering---precise specification, portability across backends, and
semantics-preserving optimization. Establishing the corresponding gains
on production workloads, and hardening the conditional sanity results
and semantic-preservation arguments into mechanically verified
implementations, is the immediate next step.

\section*{Artifact Availability}\label{artifact-availability}
\addcontentsline{toc}{section}{Artifact Availability}

The full proofs of Theorems 3.1 and 3.2 and of Theorems 3.4--3.6
(Appendices A, B, and C, referenced throughout), together with the
listings of the specified-only algorithmic components (Appendix D), the
formal sketches of the future directions (Appendix E), and the exact
artifact-regeneration command lines (Appendix F) are provided in the
supplementary material, built from this repository as
\texttt{docs/AgentISA\_supplementary.pdf}; the commands are mirrored in
the repository's top-level \texttt{REPRODUCING.md}.

All artifacts are available in the repository accompanying this
submission (anonymized for review). Reproducibility anchors:

\begin{itemize}
\tightlist
\item
  \textbf{Version}: every number in Section 5 is generated from the
  exact repository state whose build produces this PDF; the artifact's
  \texttt{git\ rev-parse\ HEAD} at submission time is recorded in
  \texttt{REPRODUCING.md}, alongside the dependency lock file and the
  Python/system versions used for the reported runs.
\item
  \textbf{Clean-environment regeneration}: from a fresh clone,
  \texttt{python\ -m\ pytest\ tests/\ -q} reproduces the full-suite
  count (861 passed, 14 skipped at run time); the Section 5.6
  files---\texttt{tests/core/isa/test\_coordination.py},
  \texttt{tests/core/instructions/test\_agent\_coordination.py}, and
  \texttt{tests/system/}---reproduce the functional-test counts (exact
  command in \texttt{REPRODUCING.md}), and the Appendix F command lines
  regenerate each results table from scratch (including the real-gateway
  runs, whose configuration---endpoint aliasing, model id, decoding
  parameters---is documented inline in the harness with a mock-backend
  fallback for environments without gateway credentials).
\item
  \textbf{Claim-to-code map}: \texttt{REPRODUCING.md} also carries a
  table mapping each paper claim to the implementing module, the test
  that exercises it, and the raw-results file that backs its numbers,
  flagging the specified-only components explicitly.
\end{itemize}

The artifact contains the AgentISA implementation
(\texttt{agentflow.core.isa}), the microbenchmark and real-LLM benchmark
suites with their raw-results files, the system-level demonstration and
four-task characterization harnesses with their tests, the real-gateway
compiler/hot-path harness with its offline-mock and real-gateway
results, and the 43 multi-agent functional tests; the full per-file
inventory is in Appendix F.5 of the supplementary material.

\end{document}
